% 来源: 19c9fbb7-1832-8871-8000-0000523133ac_5.1_虚拟网络与隧道.pdf
% 提取时间: 2026-02-28T22:51:34+08:00

5.1 虚拟网络与隧道
我们来聊一个你可能没想到的问题，结束对IP的介绍。这个问题在今天（2025年6月6日）变得越来越重
要。到目前为止，我们的讨论主要集中在如何让不同网络上的设备自由地互相通信。这其实是互联网的
核心目标——比如，每个人都希望能给任何人发邮件，一个新网站的创建者也希望尽可能多的人能访问
他们的网站。然而，在很多情况下，这种“完全开放”的连接方式并不合适，反而需要更加受控的通信方
式。虚拟专用网络（VPN）就是一个典型的例子。
“VPN”这个词被用得太多了，定义也五花八门，但我们可以简单理解：先搞清楚什么是私有网络，再来
看什么是VPN。假设一家公司有几个办公地点，它通常会租用电话公司的专线，把这些地点连接起来，
形成一个私有网络。在这个网络里，数据只能在公司内部的各个办公地点之间传输，这样可以保证安
全。要让这个私有网络“虚拟化”，就是把那些专属的专线换成一种共享的网络资源，同时仍然保持类似
的功能。比如，虚拟电路（VC）就是一个不错的选择，因为它可以在公司各个站点之间建立逻辑上的点
对点连接。
举个例子：如果公司X有一条从站点A到站点B的虚拟电路，那么它可以在站点A和站点B之间传输数据。
但公司Y的数据包是没法直接传到站点B的，除非公司Y自己也建立了一条到站点B的虚拟电路。而且，这
种虚拟电路的建立可以通过管理手段进行限制，避免公司X和公司Y之间的“串线”。
图1展示了两家公司的私有网络。在图82(b)中，这两家公司的网络都被迁移到了一个基于虚拟电路的共享
网络上。虽然它们现在共用了同样的传输设施和交换机，但各自的私有网络特性（比如只能在公司内部
传输数据）仍然保留了下来。这就是所谓的虚拟专用网络（VPN）。




 图1. 虚拟专用网络的一个示例：(a) 两个独立的专用网络；(b) 共享公共交换机的两个虚拟专用网络

                                                      1
在图1中，展示了一个基于虚拟电路网络（如ATM）的站点间受控连接架构。类似的功能亦可通过IP网络
实现，以提供等效的连接服务。然而，直接将各企业的站点接入一个统一的互联网络并非可行方案，因
为这可能导致公司X与公司Y之间形成非预期的连接，而这恰恰违背了安全隔离的设计初衷。为解决这一
问题，引入了一种新的技术概念——IP隧道。
IP隧道可被视为一种虚拟的点对点链路，逻辑上连接一对节点，而这些节点可能在物理上被多个中间网
络分隔。该虚拟链路通过在隧道入口路由器上配置远端隧道路由器的IP地址来建立。当入口路由器需要
通过该虚拟链路传输数据包时，它会将原始数据包封装到一个新的IP数据报中。此时，封装后的IP头部的
目标地址被设置为隧道远端路由器的IP地址，而源地址则为封装路由器自身的IP地址。通过这种方式，数
据能够在逻辑上隔离的路径中安全传输，从而有效避免不同企业站点之间的非授权通信。




       图2. 一个贯穿互联网的隧道。18.5.0.1 是 R2 的地址，R1 可以通过互联网访问到它
在隧道入口处路由器的转发表中，这个虚拟链路看起来与普通链路非常相似。例如，考虑图2中的网络。
已经从 R1 到 R2 配置了一个隧道，并为其分配了虚拟接口编号为 0。因此，R1 中的转发表可能如表13
所示。
NetworkNum                     NextHop
1                              Interface 0
2                              Virtual interface 0
Default                        Interface 1
路由器 R1 配备了两个物理接口和一个虚拟接口。其中，物理接口 0 连接到网络 1，而物理接口 1 连接到
一个规模较大的互联网络，并作为转发表中未匹配到更具体路由条目时的默认出口。此外，R1 的虚拟接
                                                         2
口用于隧道通信。假设 R1 从网络 1 接收到一个目标地址属于网络 2 的数据包。根据其转发表，该数据
包应通过虚拟接口 0 发送。为实现这一操作，路由器会将原始数据包封装在一个新的 IP 头中，新 IP 头
的目标地址设置为 R2（18.5.0.1），随后按照常规流程继续转发该数据包。由于 R2 的地址属于网络号
为 18 的网络（而非网络 1 或网络 2），因此该封装后的数据包将通过默认接口转发至互联网络。
一旦数据包离开 R1，它对外界而言表现为一个普通的、目标地址为 R2 的 IP 数据包，并按照常规方式
被互联网络中的路由器逐跳转发，直至抵达 R2。当 R2 接收到该数据包后，发现其目标地址与自身匹
配，于是移除外层 IP 头并解析内层数据包的有效载荷。此时，R2 识别出内层数据包的目标地址位于网
络 2 中，并以常规方式处理该数据包。由于 R2 直接连接到网络 2，它会将该数据包转发至目标网络。
图 2 展示了数据包在传输过程中封装形态的变化。
尽管 R2 在此场景中作为隧道的终点，但这并不妨碍其执行普通路由器的功能。例如，R2 可能会接收到
一些未经过隧道封装的数据包，但只要这些数据包的目标地址是其路由表中已知可达的网络，R2 将以常
规方式转发这些数据包，确保网络通信的正常运行。
您或许会思考，为何有人会不辞辛劳地构建隧道，并在数据包穿越互联网络时更改其封装方式。这一行
为背后存在多方面的原因，其中之一是安全性。通过结合加密技术，隧道能够为跨越公共网络提供一种
高度私密的通信方式。另一个原因可能与网络功能的支持有关。例如，路由器 R1 和 R2 可能支持某些在
中间网络中并不广泛部署的功能（如多播路由）。通过隧道将这些路由器连接起来，可以构建一个虚拟
网络，在该网络中，所有具备特定功能的路由器仿佛直接相连，从而实现所需功能的无缝传递。
第三个原因在于支持非 IP 协议的数据包传输。在 IP 网络上传输非 IP 协议的数据包时，只要隧道两端的
路由器能够处理这些协议，IP 隧道即可充当一条点对点的链路，允许非 IP 数据包通过。此外，隧道还提
供了一种机制，能够强制将数据包传递到特定目的地，即使其原始头部（即被封装在隧道头部内的头
部）指示它应前往其他地方。由此可见，隧道技术是一种强大且灵活的工具，可用于在互联网络中构建
虚拟链路。事实上，这项技术的通用性极高，以至于可以递归使用，最常见的应用场景便是通过 IP 隧道
传输 IP 数据包。
然而，隧道技术并非没有弊端。首先，它会增加数据包的长度，这对于短小的数据包来说可能导致显著
的带宽浪费。较长的数据包则可能因超过路径的最大传输单元（MTU）而触发分片操作，而分片本身又
会引入一系列问题，如性能下降和潜在的丢包风险。其次，隧道两端的路由器可能面临性能开销，因为
它们需要执行额外的操作，包括添加和移除隧道头部，这比普通的转发操作更加复杂。最后，负责配置
隧道并确保路由协议正确处理隧道的管理实体也需要承担一定的管理成本，包括监控、故障排查和维护
工作。因此，在实际应用中，必须权衡隧道技术带来的优势与潜在的性能和管理挑战。




                                                    3


[新增-2026-03] 虚拟机动态迁移与网络支持

虚拟机动态迁移（VMotion）是云计算数据中心的关键技术，它允许在不中断服务的情况下将运行中的虚拟机从一台物理服务器迁移到另一台。这对网络提出了特殊要求：

[新增-2026-03] VMotion的网络挑战

传统数据中心网络中资源（虚拟机）分配不能跨越3层设备，这限制了虚拟机迁移的范围。为克服这一缺陷，云计算数据中心网络采用两种解决思路：

1. 设计"大二层"网络
通过扩展VLAN或使用VXLAN等隧道技术，扩大虚拟机分配的范围，使得虚拟机可以在更大范围内迁移而不改变IP地址。

2. 借助隧道技术
使虚拟机的分配能够在"虚拟2层"中得以进行。通过MAC-in-IP或MAC-in-UDP隧道，将二层流量封装在三层网络中传输，实现跨三层设备的虚拟机迁移。

[新增-2026-03] 虚拟机迁移的协议支持

为实现快速便捷的虚拟机迁移，需要以下网络支持：

 ● 预拷贝（Pre-copy）迁移：在迁移过程中，源主机持续将内存页面复制到目标主机，同时记录被修改的页面（脏页），迭代直至脏页数量足够小，然后暂停虚拟机传输剩余脏页。这要求网络提供足够的带宽支持大量内存数据的快速传输。

 ● 后拷贝（Post-copy）迁移：先传输最小状态使虚拟机在目标主机快速启动，然后在运行过程中按需从源主机获取内存页面。这要求网络提供低延迟以保证页面获取不会显著影响性能。

 ● 实时迁移的网络优化：包括压缩传输数据、去重相同内存页面、预测热页面优先传输等技术，都依赖于网络的高带宽和低延迟特性。

[新增-2026-03] 数据中心网络与VMotion的协同设计

现代数据中心网络设计充分考虑虚拟机迁移需求：

 ● 网络虚拟化：通过软件定义网络（SDN）技术，实现虚拟机网络策略的自动跟随迁移，确保迁移后网络配置（ACL、QoS等）自动生效。

 ● 分布式虚拟交换机：在多台物理服务器上构建统一的虚拟交换机，对上层虚拟机屏蔽底层物理分布，实现透明的跨主机虚拟机通信和迁移。

 ● 东西流量优化：由于虚拟机迁移主要产生数据中心内部的东西向流量，网络拓扑设计（如Fat-Tree、Clos）专门优化东西向带宽，避免传统树形结构的汇聚层瓶颈。

\section{云时代的虚拟网络：从VPN到VPC的演进}

回望2010年代初，当我们今天讨论阿里云的VPC专有网络时，很少有人还记得阿里云最初的网络架构——Classic网络。但理解Classic时代，是我们理解整个虚拟网络演进史的关键起点。正如研究任何复杂系统一样，我们不能从最终形态入手，而必须从它的起点开始，理解那些看似"不得已而为之"的设计背后，究竟隐藏着怎样的技术权衡。

图4展示了阿里云虚拟网络从2012年到2019年的完整演进路径。这条演进路线可以分为四个阶段：\textbf{Classic时代}（2012-2015）以AVS v1/v2和MAC NAT/ARP Proxy为核心，解决了快速上线的问题；\textbf{DPDK-AVS时代}（2016-2017）通过用户态网络处理突破性能瓶颈，为高性能实例奠定基础；\textbf{VPC + HAVS时代}（2017-2018）引入Overlay网络和Session化架构，实现了虚拟网络与物理网络的解耦；\textbf{规模化与智能化时代}（2018-2019）通过分层Overlay、公网上移、IPv6双栈等技术创新，支撑了集团网络和公共云的大规模运营。

\begin{verbatim}
  阿里云虚拟网络演进路线（2012-2019）

  2012        2015          2017           2018           2019
    │          │             │              │              │
    ▼          ▼             ▼              ▼              ▼
  ┌────┐    ┌──────┐     ┌───────┐     ┌────────┐     ┌────────┐
  │Classic│ → │DPDK- │  → │VPC 1.0│ → │HAVS + │ → │规模化  │
  │网络   │   │AVS   │     │+ Overlay│   │DPDK轮转│   │+ IPv6  │
  └────┘    └──────┘     └───────┘     └────────┘     └────────┘
  │          │             │              │              │
  ├─大二层   ├─用户态转发   ├─Overlay解耦  ├─Session化   ├─分层卸载
  ├─MAC NAT │ ├─VHOST-USER ├─VXLAN/VNI  ├─热升级      ├─公网上移
  ├─ARP     │ ├─SR-IOV    ├─vRouter    ├─热迁移支持  ├─NAT64
   Proxy   │ ├─零拷贝    ├─VSwitch    ├─SmartNIC  ├─Classic迁移
            │            │            │ 卸载       │
            │            │            │            │
            └──── 带宽墙突破 ────┘            └─ 规模挑战 ──┘

       图4. 阿里云虚拟网络演进路线图。
            横轴为时间，纵向按照"问题→方案→新问题→新方案"的逻辑展开。
\end{verbatim}

\subsection{Classic时代：专线租用与隔离的代价}

2012年，阿里云上线了第一代云服务器ECS。那时候的阿里云工程师们面临一个紧迫的问题：在资源有限、研发周期紧张的情况下，如何快速实现云主机的网络功能？答案是一个朴素的方案——Classic网络。为什么要叫它"Classic"？这个名字本身就暗示了它的定位：这是最经典、最直接的网络模式，虚拟机和物理网络共享同一个二层平面，网络地址转换（NAT）、安全隔离全都靠软件层面的虚拟交换机实现。

如果我们把时间拨回到那个年代，会发现阿里云的选择其实代表了整个云计算行业的共同选择。不只是阿里云，AWS的EC2早期、腾讯云、UCloud等早期云计算厂商，几乎都采用了类似的网络架构。这不是因为大家缺乏创新意识，而是因为在那个阶段，有太多更紧迫的问题需要解决——计算虚拟化、存储虚拟化、调度系统......网络只是其中一环。Classic网络"简单直接"的特性，恰好适应了快速迭代的开发节奏。

让我们来仔细剖析Classic网络的架构特征。

\textbf{扁平的大二层架构}

Classic网络的核心特征是"扁平的大二层"。这里说的"扁平"，是指整个Classic网络在二层是完全打通的，所有虚拟机都共享同一个二层广播域。这带来一个直接的好处：虚拟机之间可以直接通过MAC地址通信，转发路径简单，延迟低。在Classic网络中，从VM A到VM B的数据包可能只需要经过：VM A的虚拟网卡 $\rightarrow$ 宿主机上的AVS $\rightarrow$ 物理交换机 $\rightarrow$ 目标宿主机AVS $\rightarrow$ VM B。整个路径最多只有3跳（如果两台虚拟机在同一台宿主机上，则只有1跳）。

然而，这种架构也埋下了一个隐患——当虚拟机数量从几十台增长到几万台时，二层广播域的规模就成了一个严峻的挑战。想象一下，如果每个虚拟机每分钟发送一次ARP请求（这在正常的网络启动过程中很常见），那么几万台虚拟机就会产生每秒数百次的ARP广播。这些广播报文会泛洪到整个大二层网络，给物理交换机带来巨大的处理压力。

更糟糕的是，大二层架构使得网络的隔离性变得脆弱。在Classic网络中，所有租户的虚拟机都在同一个二层域内，虽然有VLAN和安全组来提供隔离，但这种隔离本质上是"软"的——如果配置错误，或者虚拟机伪造了不属于它的MAC地址，就可能导致"串线"风险。

\textbf{AVS：阿里云自研虚拟交换机}

在Classic网络中，有几个关键技术组件协同工作。首先是AVS——阿里云自研的虚拟交换机（Alibaba Virtual Switch）。AVS部署在每台物理主机（Host）上，充当虚拟机的网络入口和出口。它的逻辑架构基于Linux的Netfilter框架构建，虚拟机通过虚拟网卡设备连接到AVS。

具体来说，虚拟网卡的前端连接虚拟机，后端连接到Host上的Static Bridge。Static Bridge是Linux内核提供的一种软件桥接设备，它可以将多个网络接口连接在一起，形成一个二层交换域。在AVS中，Static Bridge扮演了"内部交换机"的角色，所有连接到同一台宿主机的虚拟机，通过Static Bridge直接通信，无需经过物理网络。

虚拟机的数据通路分为TX（发送）和RX（接收）两个方向。让我们详细分析TX方向的数据流程：

第一站是Static Bridge。当虚拟机发送一个网络报文时，这个报文首先进入Static Bridge。Static Bridge根据目标MAC地址查找内部的MAC表，如果目标MAC在同一台宿主机上，就直接转发给对应的虚拟机；如果目标MAC不在本地，就将报文发送到物理网卡。

第二站是QoS模块。QoS（Quality of Service）模块负责实现网络限速功能。在云环境中，不同规格的虚拟机有不同的带宽上限（如1Mbps、100Mbps、1000Mbps等）。QoS模块通过Linux的tc（traffic control）工具，对每个虚拟机的出方向流量进行整形，确保带宽不超过规格限制。

第三站是ARP Proxy模块。这个模块是解决大二层广播风暴的关键，我们稍后会详细分析。

第四站是ACL模块。ACL（Access Control List）模块实现安全组控制——这是云平台最重要的安全机制之一。当虚拟机发送或接收报文时，ACL模块根据安全组规则判断是否允许通过。安全组规则本质上是一系列iptables规则，它们可以基于源/目标IP、端口、协议等条件对流量进行过滤。

第五站是VLAN Conf模块。虽然Classic网络是"扁平"的，但在某些场景下，仍然需要VLAN来提供额外隔离。VLAN Conf模块负责为报文打上或去除VLAN标签。

第六站是MAC NAT模块。这个模块是Classic网络的核心创新之一，我们稍后会详细分析。

第七站是Static Bridge和NIC Bonding。最后，报文经过物理网卡发出。NIC Bonding（网卡绑定）模块将多个物理网卡绑定在一起，提供更高的带宽和更好的可靠性。

\textbf{MAC NAT：地址空间压缩的艺术}

说到这里，你可能会问：为什么需要MAC NAT？这个问题触及了Classic网络的一个核心限制。

我们知道，物理交换机的MAC地址表（Neighbor表或FDB——Forwarding Database）容量是有限的。以典型的数据中心交换机为例，其MAC表容量通常在32K到128K之间。这意味着，交换机最多只能记住这么多MAC地址条目。

当虚拟机数量增长到数万台时，如果所有虚拟机都使用真实的、唯一的MAC地址，交换机根本无法存储这么多表项。更糟糕的是，MAC表耗尽后，新的MAC地址将无法学习，导致交换机将报文泛洪到所有端口——这会进一步加剧广播风暴。

因此，AVS引入了MAC NAT机制。具体来说，同一台物理主机上的多台虚拟机，在对外通信时共享同一个"公共MAC地址"（PubMAC）。当报文从虚拟机发出时，AVS将虚拟机的真实MAC地址转换为PubMAC；当报文从外部进入时，AVS再将PubMAC转换回虚拟机的真实MAC。

这种转换是如何实现的呢？答案是通过"内部MAC表"。AVS维护了一张内部MAC表，记录了每个虚拟机真实MAC与PubMAC的映射关系。这张表只存在于AVS内部，物理网络完全不可见。从物理网络的角度看，所有来自同一台宿主机的报文，都来自同一个MAC地址。

MAC NAT的效果是显著的。假设一台宿主机运行了100台虚拟机，如果没有MAC NAT，就需要交换机学习100个MAC地址；有了MAC NAT，只需要学习1个MAC地址。这将物理网络中的MAC地址数量降低了两个数量级。

\textbf{ARP Proxy：广播抑制的智慧}

另一个Classic网络必须解决的关键问题是广播风暴。由于Classic网络是二层互通的，如果大量广播、组播报文泛洪到网络上，对网络设备的压力可想而知。这些报文会消耗交换机的CPU和带宽资源，影响正常业务流量的转发。

AVS的解决方案是让虚拟机触发的广播（组播）终结在Host上，不让它们进入物理网络。但ARP报文是个例外——因为ARP是获取对端MAC地址的必要途径，而ARP本身也是一种广播报文。

为了解决这个问题，AVS中实现了ARP Proxy组件。ARP Proxy的核心思想是"代理"虚拟机的ARP请求，使得大多数ARP请求可以在本地解决，而无需泛洪到整个网络。

让我们来详细分析ARP Proxy的工作原理。

当虚拟机发出ARP请求报文（询问"IP X的MAC地址是什么？"）时，ARP Proxy首先截获这个报文，而不是让它直接发出。然后，ARP Proxy查看本地的ARP Cache——这是一个存储了最近解析结果的缓存表。如果在Cache中找到有效的条目（即IP X的MAC地址已经被解析过），ARP Proxy就按Cache条目封装ARP应答报文发回给虚拟机。整个流程到此就结束了，虚拟机不会感知到这个ARP请求曾经被"代理"过。

但如果Cache没有命中（这通常发生在第一次解析某个IP时），ARP Proxy就会"替"虚拟机发送ARP请求。它首先修改ARP请求报文中的原MAC地址为PubMAC，然后转发ARP请求报文到物理网络。当ARP Proxy收到来自交换机的ARP应答报文后，它会更新本地Cache，然后将应答报文转发给虚拟机。

ARP Proxy还有一个重要的功能：ARP流控。如果某台虚拟机短时间内发送大量ARP请求（ARP Flood），ARP Proxy可以对其进行限速，防止对整个Classic网络造成影响。这就好比是一个"门卫"，确保每个租户的ARP行为不会影响到其他租户。

\textbf{Classic网络的局限性}

Classic网络的这些设计，在早期很好地支撑了阿里云的业务发展。但随着业务规模的增长，它的局限性也越来越明显：

首先是地址空间不足的问题。Classic网络中的所有虚拟机共享一个内网地址段。在阿里云早期，这个地址段是172.16.0.0/12。随着虚拟机数量的增长，这个地址空间逐渐变得拥挤，不同租户之间可能需要通过NAT来区分，增加了复杂度。

其次是迁移域受限的问题。在Classic网络中，虚拟机的迁移范围受限于大二层网络的规模。当虚拟机需要跨机房迁移时，往往需要改变IP地址，这对应用来说是一个挑战。

再次是安全隔离不够的问题。虽然有VLAN和安全组，但大二层架构使得跨租户的流量可能经过相同的物理设备，理论上存在被窃听的风险。

最后是与物理网络强耦合的问题。Classic网络的很多行为依赖于物理网络的配置（如交换机的MAC表容量、STP协议等）。当物理网络发生变更时，可能需要同步调整虚拟网络配置，运维复杂度高。

正是这些问题的积累，推动着阿里云踏上了向VPC演进的道路。在接下来的章节中，我们将详细探讨VPC的设计理念和实现技术。

\subsection{软件虚拟交换机的诞生：DPDK-AVS的登场}

时间来到2016年，阿里云的工程师们开始着手解决一个困扰已久的问题：虚拟交换机的性能瓶颈。在Classic网络中，AVS运行在Linux内核空间，每次网络报文的处理都需要经过内核协议栈的多个模块。

让我们以一次虚拟机到虚拟机的通信为例，分析报文处理的具体路径。当虚拟机A向虚拟机B发送一个数据包时，这个数据包首先被virtio-net驱动接收。virtio-net是一个半虚拟化的网卡驱动，它允许虚拟机高效地与宿主机共享网络设备。数据包进入内核后，会经过Netfilter框架的多个hook点——这是Linux内核用于网络包过滤和修改的通用框架。

在NF_INET_PRE_ROUTING hook点，Netfilter会检查是否配置了相应的iptables规则。如果有涉及NAT、过滤等功能的规则，内核需要执行这些规则，这通常涉及大量的内存访问和条件判断。然后，数据包到达PREROUTING链，进行路由决策。如果目标是本地，就将数据包递交给上层协议栈；如果目标是其他主机，就进入FORWARD链。

经过内核协议栈的"重重关卡"后，数据包终于到达目标虚拟机的接收队列。此时，虚拟机需要处理一个中断，通知有数据到达。Linux内核的中断处理机制涉及上下文切换、锁竞争等开销。以一次中断处理为例：硬件中断触发 $\rightarrow$ CPU跳转到中断处理程序 $\rightarrow$ 关中断 $\rightarrow$ 处理中断 $\rightarrow$ 开中断 $\rightarrow$ 唤醒等待进程。这个过程本身就有几十微秒的延迟。

你可能会觉得，毫秒级的延迟对于大多数互联网应用来说应该可以接受。确实，对于网页浏览、文件下载等场景，几毫秒的额外延迟难以被用户感知。但对于那些对网络性能要求极高的业务——比如高频交易、实时游戏、视频直播——这个延迟是不可接受的。

以高频交易为例，交易指令需要在毫秒甚至微秒级别送达交易所服务器。在"快就是利润"的交易世界里，1毫秒的延迟可能意味着数十万元的损失。以视频直播为例，为了保证画面流畅，需要在50毫秒内完成数据的发送和确认。如果网络延迟超过这个阈值，观众就会感受到明显的卡顿。

更重要的是，随着云计算业务的蓬勃发展，单台物理主机上运行的虚拟机数量越来越多，网络吞吐量成为制约云服务品质的关键瓶颈。一台配置了100Gbps网卡的物理服务器，如果CPU全部用来处理网络报文，每秒可以处理约1.5亿个64字节的小包。但实际的网络处理还包括协议解析、安全检查、流量整形等功能，这些都会消耗额外的CPU资源。

问题的根源在于：传统的虚拟交换机架构，是用软件来实现网络转发的，而软件转发必然受制于CPU的处理能力。当网络带宽从10Gbps增长到25Gbps、50Gbps甚至100Gbps时，即使CPU全部用来做网络转发，也无法线速处理所有报文。这就是所谓的"带宽墙"。

\textbf{DPDK：绕过内核的数据平面}

DPDK（Data Plane Development Kit）的出现，为解决这个问题提供了一个全新的思路。DPDK是Intel发布的一套高性能数据包处理开发工具包，它的核心理念是"绕过内核"——让数据平面直接在用户态处理，从而避免内核协议栈的开销。

要理解DPDK的设计哲学，我们首先需要理解为什么传统内核网络栈的性能如此受限。

第一个原因是上下文切换。当应用程序需要发送或接收网络数据时，需要在内核态和用户态之间进行上下文切换。这种切换涉及寄存器保存/恢复、内存映射切换等操作，每次切换都有微秒级的开销。对于高频率的网络操作，这种开销是不可忽视的。

第二个原因是内存复制。在传统内核网络栈中，数据包在用户缓冲区和内核缓冲区之间需要多次复制。以一个接收到的数据包为例：网卡DMA将数据放入内核缓冲区 $\rightarrow$ 内核协议栈处理 $\rightarrow$ 复制到用户缓冲区。这至少涉及两次内存复制操作。

第三个原因是中断驱动。每当网卡收到数据包，都会触发一个硬件中断。CPU需要停止当前任务来处理这个中断，然后才能处理数据包。对于小包高吞吐量的场景，中断的频率可能非常高（如每秒数百万次），这会严重占用CPU资源。

第四个原因是锁竞争。在多核系统中，多个CPU核心可能同时访问共享的网络数据结构（如路由表、连接跟踪表等）。这需要使用锁来保证一致性，而锁的开销在高并发场景下是显著的。

DPDK通过一系列技术创新来解决这些问题：

首先是用户态驱动（UIO/VFIO）。DPDK使用用户态驱动来直接访问网卡，完全绕过了内核。当网卡收到数据包时，直接DMA到用户态缓冲区，应用程序可以直接访问，无需内核介入。

其次是大页内存（HugePages）。传统Linux使用4KB的内存页，这导致TLB（Translation Lookaside Buffer）缓存命中率低。每次内存访问都需要查询页表，页表遍历是内存访问中最耗时的操作之一。DPDK使用2MB甚至1GB的大页，减少了TLB miss的概率。

再次是无锁数据结构。DPDK为多核并行处理设计了无锁队列、环形缓冲区等数据结构，避免了锁竞争带来的开销。

最后是轮询模式驱动（PMD）。DPDK不使用中断来通知数据包到达，而是让CPU持续轮询网卡的接收队列。虽然这会增加一些CPU占用，但避免了中断处理的开销，对于高吞吐量场景是值得的。

\textbf{DPDK-AVS：阿里云的高性能虚拟交换机}

2017年，阿里云推出了基于DPDK的全新虚拟交换机——DPDK-AVS（也称为DPDK-Apsara Virtual vSwitch）。DPDK-AVS的架构设计与Classic AVS有本质的不同。

在Classic AVS中，虚拟机通过virtio-net设备和Tap设备连接到内核空间的Bridge，所有报文处理都在内核完成。这种架构的优势是简单——它利用了Linux内核已有的网络功能，不需要额外的开发工作。但缺点是性能受限，如前所述。

而在DPDK-AVS中，虚拟机通过VHOST-USER后端直接接入DPDK-AVS，报文处理完全在用户态进行。VHOST-USER是QEMU/KVM提供的一种 virtio 后端实现，它允许虚拟机与用户态的DPDK进程共享大页内存，避免了内核桥接的开销。

这种架构带来的优势是革命性的。

\textbf{收发包模式的改变}：传统virtio-net使用中断驱动的模式，当网卡的接收队列收到报文时，会触发一个中断通知虚拟机处理。这个中断处理过程本身就有几十微秒的延迟。在DPDK-AVS中，虚拟机通过VHOST-USER与DPDK-AVS共享一个大页内存区域，DPDK-AVS将收到的报文直接写入共享内存，虚拟机通过轮询来读取。这种"零拷贝"的路径大幅减少了数据复制的次数和开销。

具体来说，传统virtio-net的数据路径是：网卡DMA $\rightarrow$ 内核缓冲区（sk_buff） $\rightarrow$ 复制到用户缓冲区 $\rightarrow$ 虚拟机。而DPDK-AVS的数据路径是：网卡DMA $\rightarrow$ 用户态大页缓冲区 $\rightarrow$ 虚拟机（零拷贝）。这减少了至少一次内存复制操作。

\textbf{CPU资源的优化}：DPDK-AVS可以使用较少的CPU核心来达到更高的转发性能。测试数据显示，DPDK-AVS只需要2个超线程（2HT）就能提供比Classic AVS 8个超线程更高的转发质量和吞吐量。这意味着，DPDK-AVS可以将更多CPU资源留给用户的业务虚拟机。

以阿里云的网络增强型实例为例。该实例配备了专门的DPDK-AVS处理单元，在标准测试配置（2个Intel Xeon Platinum处理器，256GB内存，50Gbps网络）下，实测小包转发性能可达1500万pps（64字节UDP包），端到端延迟低于8微秒。如果使用Classic AVS来实现同样的性能，测试数据显示Classic AVS通常只能达到每秒100-200万包——这意味着DPDK-AVS的性能提升达到7-15倍。而DPDK-AVS只需要占用2-4个CPU超线程（2-4HT），Classic AVS要达到同等性能则需要占用8个甚至更多的超线程。从每HT产出看，DPDK-AVS的效率是Classic AVS的4倍以上。

\textbf{网络与存储流量的隔离}：在传统架构中，网络流量和存储流量共享物理网卡和CPU资源。当存储I/O负载较高时，网络性能会显著下降——因为CPU需要同时处理网络中断和存储I/O，导致"争抢"现象。

DPDK-AVS通过SR-IOV（Single Root I/O Virtualization）技术解决了这个问题。SR-IOV允许将一个物理网卡虚拟成多个VF（Virtual Function），每个VF可以独立分配给不同的虚拟机或功能模块。DPDK-AVS使用独立的VF处理网络流量，并通过DPDK独占一组CPU核心进行转发。这样，网络流量和存储流量在硬件层面就隔离了，不会相互影响。

\textbf{全用户态的可靠性}：由于所有网络处理都在用户态进行，即使DPDK-AVS出现程序bug（如内存越界、空指针解引用等），也只会导致用户态进程崩溃或重启，而不会影响整台物理服务器的稳定性。这与内核态模块的"一挂全挂"形成鲜明对比。在Classic AVS时代，一次内核模块的panic可能导致整台服务器重启，影响数千台虚拟机；而DPDK-AVS的故障影响范围被限制在进程级别。

\textbf{性能数据与业务收益}

阿里云在不到一年的时间内，新部署了250多个DPDK-AVS集群，支撑了6万多台物理机的网络增强实例。这些集群分布在阿里云的数据中心，为数百万用户提供高性能网络服务。

在第三方性能评测中，DPDK-AVS的各项指标遥遥领先于竞争对手。以小包转发性能为例，DPDK-AVS实测可达1500万pps，而Classic AVS通常只能达到100-200万pps。延迟方面，DPDK-AVS的端到端延迟实测低于8微秒，而Classic AVS通常在100微秒以上——差距达到12倍以上。

尤其是在双十一等流量高峰期，网络增强实例更是供不应求。许多电商、游戏、金融等对网络性能敏感的企业，纷纷选择使用DPDK-AVS支撑的网络增强实例来应对流量高峰。

DPDK-AVS的成功，不仅解决了当时的性能瓶颈，也为后续的虚拟网络演进奠定了技术基础。它证明了用户态网络处理的可行性，为进一步探索硬件加速打开了大门。

然而，软件虚拟交换机终究有其物理极限。随着云上业务规模的进一步扩大，阿里云开始面临新的挑战——如何在更低的成本下实现更高的性能？这推动着他们走向了硬件加速的道路。在接下来的章节中，我们将探讨AVS如何逐步引入硬件卸载能力，以及HAVS架构的设计理念。

\subsection{VPC 1.0：从二层隔离走向三层虚拟化}

如果说Classic网络代表了"虚拟机与物理网络共处一个平面"的思路，那么VPC（Virtual Private Cloud）则代表了一种完全不同的哲学——让虚拟网络与物理网络解耦，各自独立演进。

\textbf{VPC的核心思想}

VPC的核心理念是通过Overlay网络，在三层IP网络之上构建一个虚拟的二层网络。Overlay这个词本身就有"覆盖"的意思——我们可以把它理解为在现有的物理网络之上，再构建一层"虚拟网络"。这层虚拟网络有自己的地址空间、自己的拓扑结构、自己的路由规则，与物理网络完全隔离。

在VPC中，每个用户（或者说每个租户）拥有自己独立的虚拟网络。这个网络与物理网络完全隔离——物理网络只负责传输IP报文，而这些IP报文实际上是"承载着"虚拟网络报文的"隧道报文"。从物理交换机的角度看，它只看到一堆普通的IP报文；只有到了隧道的两端（VTEP——VXLAN Tunnel End Point），隧道报文才会被解封装，露出里面的虚拟网络报文。

这种设计的优势是显而易见的。

首先是独立的地址空间。在Classic网络中，所有虚拟机共享一个内网地址段，不同租户之间只能通过NAT来区分。这就好比是一个大办公室里的所有员工共用一个电话号码——虽然可以通过分机号来区分，但毕竟不够直观和灵活。而VPC允许每个用户使用自己的私有IP地址空间。以阿里云为例，VPC支持RFC 1918定义的私网地址（10.0.0.0/8、172.16.0.0/12、192.168.0.0/16），用户可以在这些范围内自由规划地址，彻底消除了地址冲突的隐患。

其次是更强的安全隔离。由于物理网络根本不认识虚拟网络的MAC地址和VLAN标签，租户之间不可能发生"串线"。即使两个租户使用了相同的私有IP地址（如10.0.1.0/24），在VPC层面它们也是完全隔离的——因为每个VPC有独立的隧道标识（如VXLAN的VNI），不同VPC的流量在物理网络上不会混淆。

再次是灵活的路由控制。VPC提供了"软件定义的路由"——在Classic网络中，虚拟机的路由行为是由物理交换机和AVS共同决定的，用户难以定制。而VPC将路由控制权完全交给用户：用户可以在VPC内创建自定义的路由表，控制虚拟机之间的流量走向；可以设置路由策略，决定哪些流量可以出VPC、哪些可以进入VPC；甚至可以通过VPN将VPC与本地数据中心连接起来，构建混合云架构。

\textbf{VPC 1.0的技术架构}

从技术实现角度来看，VPC 1.0引入了几个关键组件。

首先是vRouter（虚拟路由器）。vRouter是VPC的核心网络枢纽，负责VPC与其他网络（公网、其他VPC、本地网络）之间的流量转发。在物理网络中，路由器是连接不同网络的关键设备；在VPC中，vRouter扮演同样的角色，但它是运行在软件层面的。vRouter通常部署在VPC的网关节点上，通过软件方式实现传统的三层路由功能。

vRouter的转发能力是VPC性能的关键因素。在早期，vRouter可能面临与Classic AVS类似的性能瓶颈——随着VPC内虚拟机数量的增长，vRouter的路由表和转发性能可能成为瓶颈。这也是为什么后来阿里云在vRouter上也引入了DPDK和硬件卸载技术。

其次是VSwitch（虚拟交换机）。VSwitch是VPC中连接虚拟机的二层设备。与Classic网络中的AVS不同，VSwitch只负责VPC内部的流量交换，不直接接入物理网络。所有跨主机的虚拟机通信，都需要通过vRouter进行路由和转发。

这种"交换+路由"的分离设计，与传统物理网络中的"接入交换机+汇聚路由器"架构是一致的。区别在于，这些组件都是虚拟化的，可以根据需要进行弹性伸缩。当VPC内的虚拟机数量增长时，可以自动增加VSwitch的数量或提升vRouter的性能；当负载下降时，也可以相应缩减资源。

\textbf{VXLAN：Overlay的核心协议}

在数据层面，VPC采用了Overlay封装技术。阿里云的VPC实现使用VXLAN（Virtual Extensible LAN）协议，将虚拟网络的二层报文封装在UDP报文中进行传输。

VXLAN是VMware、Cisco等厂商在2014年提出的标准协议，它解决了传统VLAN在云计算环境下的局限性。传统VLAN只有12位的ID空间，最多支持4096个虚拟网络——对于大型云计算平台来说远远不够。VXLAN使用了24位的VNI（VXLAN Network Identifier），理论上可以支持1600万个虚拟网络。

VXLAN的工作原理如下：当虚拟机A向虚拟机B发送一个报文时，首先，虚拟机A的报文进入VSwitch。VSwitch发现目标虚拟机B在另一台宿主机上，需要跨主机通信，于是将报文封装为VXLAN格式：原始以太网帧作为负载，外面套上一层UDP头（目的端口通常是4789）、VXLAN头（包含VNI）、IP头（VTEP之间的IP）和以太网帧头。这个封装后的报文通过物理网络传输到目标宿主机，目标宿主机的VSwitch解封装后将报文交给虚拟机B。

图3展示了一个完整的VXLAN报文封装结构。从内到外，VXLAN报文依次包含：原始以太网帧（包含内层源/目标MAC、内层VLAN tag、内层IP头、内层TCP/UDP头，以及应用层数据）；VXLAN头部（8字节，包含8位Flags字段、24位Reserved字段、24位VNI字段，以及8位Reserved字段）；UDP头部（源端口由VTEP通过哈希算法生成，用于ECMP负载均衡；目的端口固定为4789，这是IANA分配的VXLAN端口）；外层IP头部（源IP为源VTEP地址，目标IP为目标VTEP地址）；以及最外层的以太网帧头（包含外层源/目标MAC，以及外层Ethernet Type）。

\begin{verbatim}
  VXLAN 报文封装结构（从内到外）
  ┌────────────────────────────────────────────────────────────────┐
  │                      外层以太网帧头 (14B)                      │
  │   外层目标MAC (6B)  │  外层源MAC (6B)  │  EtherType (2B)      │
  └────────────────────────────────────────────────────────────────┘
  ┌────────────────────────────────────────────────────────────────┐
  │                       外层IP头 (20B)                           │
  │  版本/IHL/TOS (4B) │ 总长度 (2B) │ ID/Flags/Frag (6B)        │
  │        TTL (1B)    │  协议=17 (1B) │   首部校验和 (2B)       │
  │              外层源IP地址 (4B)                                 │
  │              外层目标IP地址 (4B)                               │
  └────────────────────────────────────────────────────────────────┘
  ┌────────────────────────────────────────────────────────────────┐
  │                       UDP头 (8B)                               │
  │   源端口 (2B, VTEP哈希)  │  目的端口=4789 (2B)  │  UDP长度 (2B) │
  └────────────────────────────────────────────────────────────────┘
  ┌────────────────────────────────────────────────────────────────┐
  │                     VXLAN 头部 (8B)                           │
  │   Flags (1B, b111000)  │        Reserved (23b)               │
  │                    VNI (24b)                                   │
  │                       Reserved (8b)                             │
  └────────────────────────────────────────────────────────────────┘
  ┌────────────────────────────────────────────────────────────────┐
  │                     原始以太网帧                                │
  │  内层目标MAC (6B)  │  内层源MAC (6B)  │  VLAN/EtherType (4B)  │
  │                  内层IP头 + 传输层头 + 应用层数据               │
  └────────────────────────────────────────────────────────────────┘
       图3. VXLAN 报文封装格式。原始帧被封装在 UDP/VXLAN/IP/ETH 多层头部中，
            VNI (24位) 是区分不同虚拟网络的关键字段，VTEP 之间通过外层 IP 路由。
\end{verbatim}

VXLAN的关键创新在于"隧道端点"（VTEP）的概念。每个VTEP负责VXLAN报文的封装和解封装，它只需要知道"如何到达另一个VTEP"（这在物理网络层面通过普通路由即可实现），而不需要理解虚拟网络的拓扑细节。这种分离使得VPC可以在任何支持IP的网络上运行——即使物理网络本身不支持VXLAN，只要它能传输IP报文即可。

\textbf{为什么选择VXLAN而不是NVGRE或STT？}

在VXLAN成为云计算行业标准之前，还存在两种类似的Overlay封装协议：NVGRE（Network Virtualization using Generic Routing Encapsulation）和STT（Stateless Transport Tunneling）。

NVGRE由Microsoft等厂商提出，使用GRE（Generic Routing Encapsulation）协议来封装虚拟网络报文。与VXLAN的关键区别在于：NVGRE的负载分担字段是GRE头的"Key"字段（需要GRE解析），而不是VXLAN的UDP源端口。这意味着，在支持ECMP（等价多路径）的网络中，VXLAN更容易实现流级别的负载均衡——因为UDP源端口可以直接从内层报文头部计算出来，无需额外解析。而NVGRE的GRE Key字段分散在报文的不同位置，硬件交换机在不做深度解析的情况下难以高效地提取这个字段。

STT由Nicira（VMware）提出，也使用TCP-like的封装方式。STT的设计初衷是利用TCP的拥塞控制机制来管理Overlay流量——但这恰恰也是它的弱点：STT需要在端点模拟TCP的行为（序列号、确认号等），这增加了软件实现的复杂度；而且，STT报文在物理网络看来是一个"畸形"的TCP流（没有真正的三次握手），可能被某些防火墙或安全设备拦截。相比之下，VXLAN的UDP封装更加"干净"——UDP是一个无处不在的传输层协议，几乎所有网络设备都能正确处理它。

从硬件卸载的角度看，VXLAN也有优势。VXLAN的封装结构（UDP + 固定位置VNI）使得网卡可以更容易地实现线速封装/解封装。而NVGRE需要解析GRE头，STT需要实现完整的TCP状态机，对硬件实现的要求更高。因此，主流网卡厂商（Intel、Mellanox、Broadcom等）对VXLAN的硬件卸载支持最为完善，这也是阿里云选择VXLAN的重要原因之一。

\textbf{控制平面与数据平面的分离}

VXLAN设计的精妙之处在于：它将虚拟网络的控制平面和数据平面完全分离。数据平面（即报文的封装、转发）由VTEP执行，这些操作可以在硬件或软件中实现，取决于性能需求。而控制平面（即如何知道虚拟机B在哪个VTEP上）则可以灵活设计。

最简单的控制平面是"泛洪和学习"——当VTEP需要知道某个MAC地址在哪里时，它首先将ARP请求封装为VXLAN并泛洪到所有VTEP；收到ARP应答的VTEP记录下这个MAC与源VTEP的对应关系。这种方式简单但效率不高，对于大规模部署有局限性。

更高级的控制平面使用"集中式控制器"——如OpenFlow控制器或阿里云的飞天网络控制平面。控制器维护整个虚拟网络的状态，当虚拟机发生迁移时，控制器立即更新相关VTEP的转发表，确保流量能够正确路由。

阿里云的VPC采用了混合的控制平面策略：对于短时间的VM迁移，使用泛洪和学习机制确保可达性；对于需要快速收敛的场景，使用集中式控制器进行优化。

\textbf{VPC 1.0的意义与局限}

VPC 1.0的推出，标志着阿里云虚拟网络从"大二层"时代迈入了"Overlay"时代。这个转变的意义深远：

它解决了Classic网络面临的地址空间不足问题——每个VPC可以拥有独立的地址空间，租户之间不会冲突。

它提供了更强的隔离性——VXLAN隧道确保了不同VPC之间的流量完全隔离。

它支持更灵活的组网——用户可以自定义VPC内的网络拓扑，创建多个子网，设置路由策略。

它为混合云铺平了道路——通过VPN或专线，用户可以将VPC与本地数据中心连接起来。

然而，VPC 1.0本身也有其局限性。它仍然主要依赖软件转发，在高性能场景下可能面临性能瓶颈；它的功能相对有限（如早期的VPC不支持安全组、VPN等高级功能）；它的运维复杂度较高，需要专门的工具和培训。这些局限性推动着阿里云继续完善VPC的功能和性能，也催生了后续的硬件加速方案。

\section{硬件加速的必然：从DPDK到SmartNIC}

\subsection{AVS软件卸载：为什么软件终究不够}

2017年的双十一，创造了人类电商史上的又一个奇迹——每秒54.4万笔交易峰值，全天成交额超过1682亿元。这个数字的背后，是数以万计的云服务器在支撑。当我们把目光聚焦在这些服务器的虚拟交换机上时，会发现一个残酷的现实：即使是最优化的DPDK-AVS，也有自己的极限。

让我们来做一些具体的计算，理解软件虚拟交换机的性能边界。

假设一台物理服务器配备了2个25Gbps的网卡，总带宽是50Gbps。如果每个数据包平均大小是500字节（考虑到TCP确认包、DNS查询等小包的存在，这是一个保守的估计），那么每秒需要处理的数据包数量是：

\[
\text{每秒包数} = \frac{\text{带宽}}{\text{包大小} \times 8} = \frac{50 \times 10^9}{500 \times 8} = 12.5 \times 10^6 \text{ pps}
\]

即每秒1250万个小数据包。对于纯软件转发的DPDK-AVS来说，即使使用8个CPU核心进行并行处理，每个核心也需要处理超过150万pps的包。考虑到每个包的处理还需要执行协议解析、安全检查、流量统计等额外操作，实际的处理能力可能更低。

现代CPU每周期大约可以执行几条指令（具体取决于CPU架构和指令类型）。对于1GHz的CPU，每秒有10^9个时钟周期。如果每个包的平均处理需要200个时钟周期，那么1个CPU核心每秒可以处理约500万pps。这意味着，要达到1250万pps的吞吐量，至少需要3个满载的CPU核心。

更重要的是，这个极限还会随着网络功能的增加而降低。在实际生产环境中，虚拟交换机不仅要做基本的报文转发，还要执行复杂的功能：

首先是安全检查（ACL/安全组）。每当一个报文经过虚拟交换机时，需要与一系列ACL规则进行匹配。这些规则可以基于源/目标IP、端口、协议、MAC地址等多个字段。平均而言，每个报文可能需要匹配10-50条规则。ACL匹配的时间复杂度是O(n)，其中n是规则数量。

其次是流量整形（QoS）。QoS需要对每个虚拟机的流量进行计量和限速，确保不超过其规格限制。这涉及到令牌桶算法、队列调度等复杂操作。

再次是会话跟踪（Session/连接跟踪）。对于有状态的流量转发（如NAT、负载均衡），虚拟交换机需要维护连接状态表。每当一个报文到达时，需要查找对应的会话；如果是一个新连接，则需要创建会话条目。

最后是隧道封装/解封装。对于VXLAN流量，需要添加或去除额外的头部，这在软件中也是不小的开销。

每一项功能都会消耗CPU cycles，最终导致整体吞吐量下降。实测数据表明，当同时开启ACL、QoS、Session等功能时，DPDK-AVS的吞吐量可能下降30\%到50\%。

\textbf{成本问题：被"浪费"的CPU核心}

除了性能瓶颈，还有成本问题需要考虑。

我们知道，DPDK-AVS需要独占一组CPU核心来进行网络转发。这些核心不能运行用户的业务虚拟机，也不能运行其他系统进程。以阿里云的网络增强型实例为例，该实例规格配备了8个vCPU，其中2个（25\%）被DPDK-AVS占用，剩余6个（75\%）可以分配给用户。

这个比例看起来还可以接受。但当我们考虑整个阿里云的规模时，问题就变得严重了。假设阿里云有10万台物理服务器，每台服务器平均有4个DPDK-AVS核心被占用（用于不同类型的网络功能），那么总计有40万个CPU核心被用于网络转发。

如果这些核心可以用于运行用户虚拟机，按照每核心每年约100美元的云服务价格（这是保守估计），每年"浪费"的成本就是4000万美元。这还不包括这些核心的电力消耗、散热成本等。

因此，从成本角度来看，也迫切需要将更多的网络功能卸载到专用硬件上。

\textbf{SmartNIC：智能网卡的崛起}

正是这些挑战——性能瓶颈和成本压力——推动着阿里云将目光投向了硬件加速，特别是SmartNIC（智能网卡）。

SmartNIC是一种内置了可编程芯片的网卡，它可以在网卡上直接执行网络功能，释放CPU的计算压力。与传统的网卡相比，SmartNIC不仅仅是"更快的网卡"，它实际上是一台小型的专用计算机。

SmartNIC的核心组件通常包括：

首先是网络处理单元（NPU或网络处理器）。这是一种专门设计用于网络包处理的芯片，可以高效地执行协议解析、头部修改、校验和计算等操作。

其次是可编程逻辑（FPGA）。FPGA允许用户自定义硬件逻辑，实现特定的协议处理或算法。对于云计算场景，FPGA可以用于实现自定义的VXLAN封装、安全检查等功能。

再次是高速内存（TCAM、DDR等）。TCAM（三态内容寻址存储器）是一种特殊的内存，可以同时在多个位置进行并行查找，常用于ACL规则的快速匹配。

最后是高速接口（PCIe Gen3/Gen4、100G Ethernet等）。SmartNIC需要高速接口与服务器CPU和其他设备通信。

与传统的网卡相比，SmartNIC的核心优势在于：

\textbf{线速处理能力}：SmartNIC上的专用网络处理器可以在硬件层面线速处理网络报文，不受CPU性能的限制。即使是100Gbps的网络带宽，SmartNIC也能轻松应对，没有软件转发那种"带宽墙"的限制。

\textbf{可编程性}：现代SmartNIC采用FPGA或ASIC芯片，支持用户自定义的网络功能。这意味着，云厂商可以根据自己的需求，在网卡上实现特定的协议处理、安全检查、流量分析等功能，而不受网卡固件的限制。

\textbf{CPU卸载}：将网络功能卸载到网卡上后，服务器CPU可以完全专注于运行用户业务。据统计，一张好的SmartNIC可以"释放"相当于8-16个CPU核心的计算能力。这些释放的CPU资源可以用于运行更多的虚拟机，或者降低服务器的配置要求。

阿里云的硬件加速之路始于2017年。最早期的方案是将AVS的部分功能（如VXLAN封装/解封装）卸载到网卡上执行，而控制平面仍然运行在服务器CPU上。这种"部分卸载"的策略，降低了技术风险，也证明了硬件加速的可行性。

然而，硬件加速也带来了新的挑战。首先是架构复杂度的增加——当网络功能分布在软件和硬件两个层面时，如何保证它们之间的协调一致？其次是升级维护的困难——软件版本可以快速迭代，而硬件的升级周期通常更长。再次是成本问题——SmartNIC的价格是普通网卡的数倍，需要权衡投入产出比。

这些挑战的应对，推动着阿里云设计出了更加完善的HAVS架构。

\subsection{HAVS：可编程硬件上的虚拟交换}

回顾AVS的演进历史，我们可以看到一条清晰的技术演进脉络。

AVS v1和v2版本，虽然实现了基本的虚拟交换机功能，但架构上存在一些历史遗留问题。这些问题不是"错误"，而是"技术债务"——在快速迭代的开发阶段，这些设计可能是合理的选择；但随着系统的发展，这些设计逐渐成为瓶颈。

\textbf{对Linux的强依赖}

AVS v1/v2实际上是由多个Linux内核层面的驱动模块组成的，通过内核的Netfilter串联起来。这种设计导致对内核版本、iptables、tc等有严重依赖。

具体来说，每个功能模块（ARP Proxy、ACL、QoS等）都是作为Linux内核模块实现的，它们通过Netfilter的hook机制串联在一起。这种架构的优点是可以利用Linux内核已有的基础设施（如iptables、tc），加快开发速度。但缺点也很明显：

第一，对内核版本的强依赖。当Linux内核版本升级时，可能导致AVS模块不兼容。这限制了内核升级的频率，也增加了测试的工作量。

第二，内存分配的问题。Linux内核部分功能涉及到kmalloc等连续物理内存申请的问题。kmalloc只能分配物理连续的内存，当系统运行一段时间后，内存碎片化会导致分配失败。线上环境曾多次出现tc qdisc高阶内存分配失败导致的系统负载上升、不稳定，严重时甚至导致服务器hang。

第三，iptables的扩展性问题。ACL功能依赖于iptables，而iptables随着规则数的增加性能急剧下降。当安全组规则数量达到数万条时，iptables的匹配时间可能从微秒级增长到毫秒级。

\textbf{架构的离散性}

不同的功能模块（iptables、tc、ARP Proxy等）各自独立运行，无法共享公共的库和数据结构。这导致同一个功能模块在不同地方需要实现多遍，开发效率低下，维护成本高企。

例如，MAC地址解析功能在ARP Proxy和MAC NAT两个模块中都有实现。虽然它们的逻辑类似，但由于运行在不同的上下文（不同模块、不同内存空间），无法直接复用代码。这不仅增加了开发工作量，也增加了维护成本——当需要修改MAC解析逻辑时，需要同时修改两处。

\textbf{性能瓶颈}

由于各功能模块的架构是离散串行的，随着新业务的增加，总体性能线性下降。例如：

iptables随着规则数的增加性能急剧下降——这是因为iptables使用链表存储规则，每条规则的匹配都需要遍历链表。

QoS依赖ip-route2的tc，有巨大的锁开销——tc的队列管理涉及大量的锁操作，在多核环境下成为瓶颈。

Session管理的效率低下——v1/v2版本的AVS使用"每报文匹配"的架构，每个报文都需要遍历所有规则才能确定如何处理。

\textbf{可维护性差}

AVS v2的RPM安装包就有十几个之多。Classic和VPC功能封装在不同的安装包中，不同的虚拟化平台（Xen和KVM）上也有不同的安装包。每个模块都有自己独立的一套运维工具和脚本。

这种分散的架构带来了严重的运维挑战：升级时需要同时升级多个包，每个包之间可能存在依赖关系；回滚时需要同时回滚多个包，如果某个包回滚失败可能导致系统不一致；故障排查时需要检查多个模块的日志和状态，每个模块的日志格式和诊断命令都不一样。

由于模块之间业务上耦合但部署上独立，运维部署升级和回滚都十分复杂。一次看似简单的安全补丁，可能需要在数百台服务器上手动操作数小时。

\textbf{HAVS的架构重构}

HAVS（High-performance AVS）的设计目标，就是从根本上解决这些问题。

HAVS将所有功能统一整合，重新架构，可同时支持Classic网络和VPC，并针对硬件加速进行了优化设计。与其说HAVS是一个新版本的AVS，不如说它是一个全新的系统——虽然名字还叫AVS，但内部架构已经完全不同。

在架构层面，HAVS采用了"控制平面+数据平面"分离的设计。这种设计模式借鉴了软件定义网络（SDN）的思想，但在实现上更加务实。

控制平面负责策略下发、状态管理等逻辑，运行在服务器CPU上。它包括：

路由控制模块——管理VPC的路由表，决定报文的下一跳；
安全策略模块——管理ACL规则和安全组，负责安全检查；
会话管理模块——维护TCP/NAT会话状态，支持有状态的转发；
设备管理模块——与OpenStack/Kubernetes等云平台对接，响应资源创建/删除请求。

数据平面负责实际的报文转发，既可以在CPU上运行（软件FastPath），也可以卸载到SmartNIC上运行（硬件FastPath）。这种设计使得HAVS可以灵活地在软件和硬件之间切换，根据业务需求选择最优的转发路径。

\textbf{Session化的转发架构}

在技术实现上，HAVS引入了几个关键创新。

首先是Session化的转发架构。与传统的每报文匹配规则不同，HAVS将网络流量抽象为"会话（Session）"。

什么是会话？简单来说，会话是一系列具有相同特征的报文组成的数据流。例如，一个TCP连接就是一个会话——第一个SYN报文会创建一个会话条目，后续属于这个连接的报文都可以通过查找会话表来快速处理。

Session化的优势是显著的。对于属于同一会话的后续报文，直接根据会话表进行转发，无需重复匹配复杂的ACL规则。这大幅提升了转发效率，也减少了规则数量对性能的影响。

以一个典型的HTTP连接为例：连接建立时，第一个报文需要经过完整的处理流程——检查源IP是否被禁止、目标端口是否开放、应用层协议是否匹配等，可能涉及数十次规则匹配。但连接建立后，后续数百甚至数千个报文只需要查找会话表——这只需要一次内存访问。

会话表的存储也有讲究。HAVS使用高性能的哈希表来存储会话条目，支持O(1)时间复杂度的查找。对于有状态的转发（如NAT），会话条目还包含地址转换的信息。

会话超时管理也是一个重要的考虑。当一个会话长时间没有流量时，需要将其从会话表中删除，释放内存。HAVS使用多级超时机制：高频会话使用较短的超时，低频会话使用较长的超时。

\textbf{热升级能力}

HAVS的控制平面支持热升级——在升级过程中，新的控制平面进程可以接管旧进程管理的连接状态，实现用户无感知的平滑升级。

热升级的工作原理如下：当需要升级HAVS时，控制系统首先在目标服务器上启动新版本的进程；新进程与旧进程建立通信，接收旧进程的会话状态和配置信息；新进程验证接收到的状态后，向控制系统报告"就绪"；控制系统开始将流量切换到新进程；旧进程在完成流量切换后退出。

整个过程中，用户的业务流量不会中断。测试数据显示，对于一个包含10万会话的HAVS实例，热升级过程中的流量中断时间可以控制在100毫秒以内——这对大多数应用来说都是不可感知的。

这解决了传统AVS升级需要中断业务的问题。在v1/v2时代，升级AVS意味着需要重启虚拟机或迁移虚拟机到其他服务器，这对用户来说是有感的。

\textbf{对热迁移的支持}

虚拟机迁移是云计算的核心能力之一。在虚拟机迁移过程中，需要保持网络会话的连续性。如果虚拟机迁移后，之前建立的TCP连接全部断开，用户体验将大打折扣。

HAVS的Session架构天然支持会话状态的迁移。当虚拟机迁移到新的物理服务器时，原服务器的会话状态可以"迁移"到新服务器。

迁移过程如下：当虚拟机开始迁移时，源服务器上的HAVS将会话状态序列化并传输到目标服务器；目标服务器上的HAVS接收并解析会话状态，重建会话表；当虚拟机在目标服务器上启动后，流量会被引导到新服务器；由于会话表已经就绪，用户流量无需重新建立连接。

这种方法确保了迁移过程中连接的连续性。对于SSH会话、数据库连接等长连接应用，这种能力尤为重要。

\textbf{硬件卸载的深度优化}

HAVS不仅支持将基础的封装/解封装操作卸载到SmartNIC，还支持将安全检查、流量统计等功能的全部或部分卸载到网卡上。

具体来说，HAVS的硬件卸载分为几个层次：

第一层是基础卸载——将VXLAN封装/解封装卸载到网卡。这种卸载已经被广泛采用，它可以将隧道处理的开销从CPU转移到网卡。

第二层是安全卸载——将ACL规则的匹配卸载到网卡。网卡使用TCAM来存储规则，可以线速匹配每个报文。当规则数量较少时（通常小于1万条），这种卸载可以完全消除ACL的性能影响。

第三层是计量卸载——将流量统计和限速卸载到网卡。网卡使用专门的计数器来记录每个流的流量，使用令牌桶算法来进行限速。这可以将QoS的处理从CPU卸载，释放更多的CPU资源。

这种深度的硬件卸载，可以将SmartNIC的性能优势发挥到极致。

测试数据显示，相较于DPDK-AVS，HAVS在吞吐量、延迟、CPU利用率等指标上都有显著提升：

吞吐量提升30\%到50\%——得益于Session架构和硬件卸载；
延迟降低50\%——减少了很多不必要的数据处理；
CPU占用降低40\%——更多的处理被卸载到网卡。

特别是对于小包场景，HAVS的性能提升尤为明显——这得益于Session架构减少了规则匹配的次数，以及硬件卸载避免了CPU参与每个报文的处理过程。

HAVS的成功，不仅解决了AVS的历史遗留问题，更为阿里云的虚拟网络架构演进奠定了坚实的基础。在此基础上，阿里云进一步实现了HAVS与DPDK-AVS的协同轮转，使得线上服务器可以根据业务需求灵活切换转发模式。

\subsection{HAVS与DPDK-AVS的共存与轮转}

2018年，阿里云的虚拟交换机迎来了一个重要的里程碑——DPDK-AVS全面上线。然而，此时线上还运行着大量的HAVS服务器，这些服务器主要支撑存量用户的Classic网络实例。如果能够将这些HAVS服务器升级为DPDK-AVS，将直接获得显著收益。

\textbf{轮转的价值}

将HAVS升级到DPDK-AVS，可以带来多方面的收益：

第一，HAVS共享集群可以直接升级到网络增强集群，大幅提升资源的投入产出比。网络增强实例的单价更高，可以为阿里云带来更多收入。

第二，由于DPDK-AVS 2HT就可以提供比HAVS 8HT更高的服务质量，项目成功实施后将节省大量CPU核心。以阿里云10万台服务器的规模计算，每年可节省的成本达到数亿元。

第三，DPDK-AVS在物理网卡层面将网络流量与存储流量隔离，使用独占的CPU，彻底解决了困扰ECS多年的网络存储争抢问题。用户体验将显著提升。

第四，全用户态改造，根除了因为程序bug导致整服务器宕机的风险。这将大幅减少因服务器宕机导致的客户投诉和赔偿。

第五，统一AVS架构，避免两套架构重复开发测试、运维，显著提高开发效率，简化库存管理。维护两套架构需要双倍的开发人员和运维人员。

这，就是HAVS轮转DPDK-AVS项目的背景。

\textbf{挑战与方案分析}

然而，改造过程面临一个基本约束：不能影响用户。这意味着，我们需要一种在用户无感知的情况下，将HAVS升级为DPDK-AVS的方案。

团队最先考虑的方案是热升级——类似于HAVS自身的热升级能力，将VM和物理网卡侧的网络流量从老的AVS切换到新的AVS。

热升级AVS，就是将VM和物理网卡的流量从老的AVS切换到新的AVS。由于AVS还会基于Session提供状态ACL、SLB等业务功能，所以热升级时还涉及拷贝Session以及配置信息同步。

总结起来，热升级AVS，本质上是完成以下工作：

配置信息同步——将老AVS的配置复制到新AVS；
Session拷贝——将老AVS的会话表复制到新AVS；
VM侧网络流量切换——将虚拟机的流量从老AVS引导到新AVS；
物理网卡侧网络流量切换——将物理网卡的流量从老AVS引导到新AVS。

然而，进一步分析发现，这个方案面临几个难以克服的挑战。

首先是VM侧接入方式的差异。DPDK-AVS为了提高开发跟运维效率，运行在用户态，VM通过VHOST-USER后端直接接入DPDK-AVS。而在HAVS上，VM是通过VHOST-NET关联Tap设备，再通过Bridge将Tap设备接入HAVS。

这两种接入方式的区别是本质的。VHOST-USER使用共享内存的方式让VM和DPDK-AVS直接通信，不需要经过内核；而VHOST-NET/Tap/Bridge的方式需要经过内核协议栈。由于VM接入AVS的方式是在VM启动时指定的，无法热修改。

其次是大页内存的问题。DPDK-AVS为了提高TLB命中率，VM及VHOST-USER都使用了2M大页。而HAVS使用透明大页。VM流量切换到DPDK-AVS需要重新分配大页内存给VM，不重启VM似乎也无法做到。

再次是物理网卡侧网络流量切换的难点。由于DPDK-AVS开启了网卡VF功能，从物理网卡层面隔离网络流量和存储流量，所以升级到DPDK-AVS也需要UP/DOWN网卡，重新加载驱动。如果NC流量小于10G（单网卡容量），可以采取逐网卡切换的方式升级；但当NC流量大于10G时，就无法做到无损切换了。

综上所述，似乎没有方案能不重启VM跟网卡，将NC升级到DPDK-AVS。

\textbf{轮转热迁移的创新}

既然重启不可避免，有没有办法重启VM，但对用户没有影响呢？答案就是轮转热迁移。

轮转热迁移的思路很巧妙：先将VM热迁移到其他服务器，腾空HOST；然后在腾空的HOST上直接部署DPDK-AVS；最后再将VM迁移回来。

在这个过程中，用户只经历了VM热迁移的短暂中断（通常在毫秒级），而没有服务中断。

这个方案涉及两个关键技术。

第一项关键技术是对于各种业务热迁移都能做到毫秒级断网。热迁移对网络来说，本质是如何快速地将"VM从NC_A迁移到NC_B"这个事件通知到所有相关的网络设备。

毫秒级断网的关键在于热迁移Relay机制。当VM从源主机迁移到目标主机时，与VM保持通信的对端可能感知到MAC地址的变化（如果迁移跨网段）或者短暂的不可达（ARP缓存过期）。

Relay机制通过在迁移过程中保持一个"中转"路径，确保即使VM的物理位置发生变化，通信对端仍然能够找到VM。Relay设备会"替"VM临时响应ARP请求，直到VM迁移完成并完成新位置的宣告。

Relay的工作原理可以分为以下几个阶段：

第一阶段（迁移前）：源主机上的HAVS向Relay设备同步该VM的所有会话状态和路由信息。Relay设备建立对应的"影子"转发规则，记录该VM的原始MAC/IP和当前位置。

第二阶段（停机拷贝期间）：VM在源主机上暂停运行，最后一次增量内存被传输到目标主机。此时，VM的所有网络连接处于"冻结"状态——从对端看来，连接仍然存在（因为对端的TCP状态机尚未感知到连接中断），但没有新数据发送。

第三阶段（切换时刻）：迁移控制系统通知Relay设备"VM即将切换位置"。Relay设备将对应的"影子"规则激活，开始代为转发目的地址为该VM的所有报文。这些报文被Relay截获后，重新路由到VM的新位置（目标主机）。这个过程在毫秒级完成。

第四阶段（VM在新位置启动）：VM在目标主机上启动，HAVS立即宣告VM的新位置（通过发送无偿ARP或GARP报文）。网络中的交换机更新CAM/ARP表，将MAC/IP映射指向新的物理位置。同时，HAVS向Relay设备同步最新的会话状态，Relay逐渐将转发权交还给正常的网络路径。

第五阶段（清理）：当Relay确认所有对端已经完成路由更新后，删除本地的"影子"规则。整个Relay机制退出。

整个过程中，通信对端的TCP连接从未断开——对端的TCP状态机认为连接始终保持，而实际的数据路径在Relay的帮助下悄悄切换到了新位置。这正是"热迁移对用户无感知"的技术基础。

第二项关键技术是支持跨AVS版本的热迁移。HAVS和DPDK-AVS虽然转发架构不同，但都支持基于Session的转发模式。

在迁移过程中，HAVS可以将VM相关的Session信息同步给DPDK-AVS，使得新创建的DPDK-AVS能够立即处理该VM的流量，无需重新建立所有连接。这项技术涉及会话状态的序列化和传输，需要确保状态的完整性和一致性。

通过HAVS轮转DPDK-AVS项目，阿里云在短短几个月内完成了数万物理服务器的升级改造。数万虚拟机被成功迁移，迁移成功率达到99.99\%以上。

更重要的是，这个项目的成功经验，为后续的架构演进提供了一个可复用的模式——通过热迁移实现无感知的基础设施升级。这个模式后来被应用到更多的场景中，如内核升级、固件更新、安全补丁等。

\section{VPC的规模化挑战}

\subsection{VM热迁移：网络一致性的极限挑战}

虚拟机热迁移（Live Migration）是云计算数据中心最核心的技术之一。简单来说，热迁移允许将一台运行中的虚拟机从一台物理服务器迁移到另一台，整个过程对用户是透明的——用户的应用不会感知到底层基础设施的变化。

在前面讨论HAVS架构时（第4.2节），我们已经介绍了HAVS的Session化设计如何支持热迁移——会话状态可以随虚拟机一起迁移，确保TCP连接不会中断。本节我们将聚焦VPC Overlay网络给热迁移带来的额外挑战：VTEP更新、跨可用区双写等HAVS架构本身无法解决、VPC特有的一致性问题。

这项技术的价值在于多个方面。

首先，它实现了硬件故障的"零容忍"。当某台物理服务器出现硬件故障时——如内存错误、硬盘损坏、电源故障——运行在其上的虚拟机可以自动热迁移到健康的服务器上，用户服务不会中断。这大幅提升了云平台的可靠性指标。

其次，热迁移使得资源调度更加灵活。数据中心可以根据负载情况，将虚拟机迁移到资源更充裕的服务器上。例如，当某台服务器CPU负载过高时，可以将其上的部分虚拟机迁移到负载较低的服务器，实现负载均衡。在维护窗口期，运维人员可以先将虚拟机迁移走，对服务器进行维护，然后再将虚拟机迁回来。

再次，热迁移支持弹性伸缩。在弹性伸缩场景中，当业务负载增加时，可以快速启动新虚拟机；当负载下降时，可以将虚拟机迁移到少数服务器上，然后释放空闲服务器。

\textbf{热迁移对网络的挑战}

然而，虚拟机热迁移对网络系统提出了极高的要求。

想象一下：当你正在观看一段视频时，视频数据正通过某个网络路径传输。如果此时虚拟机突然"跳"到了另一个物理位置，而网络没有及时更新转发表，你会看到什么？答案是一段视频卡顿或者马赛克。

这个问题的根源在于：网络的状态（路由表、ARP缓存、连接跟踪表等）都是与物理位置相关的。当虚拟机迁移到新位置时，这些状态需要"跟过去"，否则就会导致通信中断。

让我们来详细分析热迁移的网络挑战。

在经典的大二层网络中，虚拟机的MAC地址与所在的二层网络直接相关。当虚拟机从一个物理服务器迁移到另一个时，如果两台服务器在同一个二层网络内（它们连接到同一个物理交换机，或者通过Trunk链路相连），MAC地址仍然有效，转发表只需要更新虚拟机的新位置即可。然而，如果两台服务器不在同一个二层网络内（比如跨交换机、跨机房），MAC地址就会失效，需要重新进行地址解析。

地址解析的过程可能需要数百毫秒到数秒不等——这对于实时应用来说是不可接受的。

\textbf{VPC Overlay的优雅解法}

VPC的Overlay网络为这个问题提供了一个优雅的解决方案。

在VPC中，虚拟机的MAC地址和IP地址都是虚拟的，与物理网络解耦。虚拟机迁移时，只需要更新Overlay网络的封装信息（比如VXLAN的VTEP——VXLAN Tunnel End Point地址），物理网络完全不需要感知。

具体来说，当虚拟机从服务器A迁移到服务器B时，VPC控制平面会立即更新VXLAN的映射表——原来指向服务器A的VTEP，现在指向服务器B。之后发往这个虚拟机的流量，会自动被路由到新的服务器。

这种"映射更新"的过程非常快速——在阿里云的实现中，控制平面可以在50毫秒内完成更新。这意味着，迁移后的短暂中断时间被控制在毫秒级别，对于大多数应用来说都是不可感知的。

但这并不意味着VPC的热迁移就没有挑战了。在实际实现中，还有几个关键问题需要解决。

\textbf{Session一致性问题}

首先是Session一致性问题。对于有状态的业务（如TCP连接、NAT会话），当虚拟机迁移到新位置时，这些会话状态需要保持连续性。

想象一下：你正在使用一个基于TCP的远程桌面连接到虚拟机。如果VM迁移后，TCP连接全部断开，你会看到桌面瞬间黑屏，然后需要重新连接。虽然重新连接可能只需要几秒钟，但这对于用户体验来说是明显的"断层"。

解决这个问题的方法是"有状态热迁移"。顾名思义，有状态热迁移不仅迁移虚拟机的内存和CPU状态，还迁移与网络相关的会话状态。

阿里云的实现方案中，会话状态作为虚拟机内存的一部分，通过热迁移的"预拷贝"（Pre-copy）机制传输到新服务器。在预拷贝阶段，源服务器先传输虚拟机的全部内存到目标服务器；然后，只传输内存中"脏"（被修改）页面的增量；这个过程反复进行，直到脏页数量足够少，才进行最终的停机拷贝。

在这个过程中，HAVS的Session状态会被同步传输。当迁移完成、虚拟机在目标服务器启动后，新服务器上的HAVS会继承原服务器的Session表项，确保用户流量不会中断。

当然，Session同步也有其局限性。如果会话数量非常大（如数十万个），同步这些会话可能需要较长时间。在极端情况下，可能需要在迁移前将一些"冷"会话丢弃，由应用程序在需要时重新建立。

\textbf{网络策略跟随的问题}

其次是网络策略跟随的问题。在云计算环境中，虚拟机的网络策略（如ACL安全组、带宽限制）通常与虚拟机本身绑定，而非与物理位置绑定。

以安全组为例，安全组定义了允许哪些流量进出虚拟机。当虚拟机迁移时，这些安全组规则需要自动"跟随"到新服务器。

阿里云的解决方案是通过集中式的控制平面来管理网络策略。当虚拟机迁移时，控制平面会立即将相关的策略下发到新的AVS/HAVS上。由于策略的下发是独立于数据平面的，所以可以实现毫秒级的策略生效。

这种设计的另一个好处是：即使虚拟机还没有完全迁移到新服务器，新服务器上的HAVS也已经配置好了相关的策略。当虚拟机开始接收流量时，安全策略已经就绪。

\textbf{迁移中断时间的优化}

再次是迁移中断时间的问题。即使有上述优化，虚拟机热迁移本身也会带来短暂的网络中断。

这个中断时间主要用于：源服务器停止发送流量、传输最后的增量数据、更新目标服务器的转发表、通知通信对端新的位置等。

阿里云通过多种技术来优化中断时间：

优化迁移协议——减少数据传输的轮次；
压缩传输数据——减少需要传输的数据量；
预测热页面——优先传输可能被访问的内存页；
并行传输——利用多个网络路径同时传输。

实测数据显示，对于典型的Web应用虚拟机（4GB内存，网络连接数小于1万），阿里云可以将迁移中断时间控制在50毫秒以内。对于大多数应用来说，这个级别的中断是不可感知的——它比一次正常的TCP超时重传还要短。

\textbf{跨可用区迁移}

最后是跨可用区迁移的挑战。在阿里云的架构中，一个地域（Region）包含多个可用区（Zone）。每个可用区是独立的物理数据中心，有自己的网络、供电、冷却系统。跨可用区迁移用于实现更高可用性——即使整个可用区发生灾难性故障，业务也可以在另一个可用区继续运行。

然而，跨可用区的虚拟机迁移涉及更长的网络路径和更多的网络设备，网络一致性的维护更加复杂。

首先，跨可用区的网络延迟更高。同一可用区内的迁移，延迟通常在亚毫秒级；而跨可用区的延迟可能在1-10毫秒。这对于热迁移的停机切换时间有直接影响。

其次，跨可用区涉及更多的网络设备。虚拟机可能需要经过多个路由器和交换机才能到达目的地，任何一个设备的故障都可能导致迁移失败。

阿里云通过"双写"机制来解决跨可用区迁移的一致性问题。在迁移过程中，虚拟机的状态同时在源和目标两个位置保持一致，直到确认目标位置已经准备好接收所有流量后，才停止源位置的写入。这种方式可以确保跨可用区迁移的强一致性，但也带来了更高的复杂度和成本。

VM热迁移技术的持续演进，体现了云计算平台对"零中断"服务的不懈追求。随着网络技术的进步和架构的优化，热迁移的适用范围不断扩大，为云计算的弹性运维提供了坚实的基础。

\subsection{集团Overlay落地：规模化的工程实践}

2018年对于阿里云来说是具有里程碑意义的一年。在这一年的双十一中，阿里云的技术架构经历了严峻的考验——单日交易峰值达到2135亿元，比上一年增长27\%。支撑这个天文数字的，是阿里集团数十万台物理服务器、数百万个虚拟机，以及覆盖全球的数据中心网络。

然而，随着业务的快速发展，集团网络面临着前所未有的挑战。

最核心的问题是：如何让数量庞大的虚拟机实例能够像在一个本地网络中一样自由通信，同时又要保证隔离性和安全性？

\textbf{双轨并行的历史架构}

在Overlay网络技术成熟之前，阿里集团采用了两种并行的网络架构。

第一种是经典的VLAN大二层网络。虚拟机可以在大二层内自由迁移，延迟低、配置简单。这种架构与阿里云公共云的Classic网络类似，所有虚拟机共享同一个二层广播域。VLAN提供了基础的隔离能力，但VLAN ID只有4096个，无法支持大量的隔离域。

第二种是基于硬件的MPLS VPN。通过路由器和交换机的硬件转发实现多租户隔离。这种架构提供了更强的隔离性和可扩展性，但配置复杂、灵活性差。每增加一个新的"租户"（业务部门），都需要在路由器上配置一系列的VRF（Virtual Routing and Forwarding）、LDP（Label Distribution Protocol）等协议。

这两种架构并行运行了多年，各自服务不同的业务场景。随着业务规模的增长，运维成本也在不断上升——需要维护两套独立的网络设备和配置体系，任何变更都需要在两个系统中同步进行。

\textbf{集团Overlay的技术挑战}

阿里云的解决方案是：将VPC Overlay技术从公共云扩展到整个集团网络。

这并不是一个简单的"复制粘贴"过程。公共云的VPC与集团内部网络有本质的差异。公共云的租户是外部客户，每个客户有自己的VPC；网络策略由平台统一管理；租户之间的隔离是强制的。而集团内部网络的"租户"是各个业务部门，它们之间可能有合作关系，希望能够方便地互相访问；网络策略可能由业务部门自行定义；隔离可以是灵活的。

将这些差异简单地罗列出来，我们就得到了集团Overlay落地需要解决的几个关键挑战。

第一个挑战是地址空间的统一规划。

不同的业务部门可能使用了相同的私有IP地址段。例如，电商部门可能使用了10.0.1.0/24作为内部网络的地址段，而云计算部门也可能使用了相同的地址段。在Classic网络时代，由于大二层的隔离，这两个部门实际上不会发生冲突。但在VPC Overlay中，如果要将这些业务接入统一的Overlay网络，就需要进行地址规划，避免Overlay隧道两端的地址冲突。

解决的方法是采用"租户ID+原始地址"的编码方案。每个业务部门被分配一个唯一的租户ID（这在阿里集团内部称为"BizID"），Overlay隧道报文中的内层地址包含租户ID信息。这样，即使两个业务部门使用了相同的内网地址（如10.0.1.5），在Overlay层面也能被正确区分——它们分别属于不同的租户。

具体来说，阿里云在VXLAN头部的VNID字段中嵌入了租户ID。VNID是24位字段，阿里云对其做了分层编码：高8位用于表示Region（地域），中间8位用于表示BizID（业务部门ID），低8位用于表示部门内的子网ID。这种编码方式使得：
\begin{itemize}
    \item 同一个业务部门在不同Region的流量，可以根据VNID高8位区分，避免跨Region的路由混淆；
    \item 同一个Region内，不同业务部门的流量根据中间8位BizID区分，确保即使地址重叠也不会串线；
    \item 同一个业务部门内部，不同子网的流量根据低8位子网ID区分，实现精细化的网络隔离。
\end{itemize}
这种编码方案的另一个好处是：VTEP可以根据VNID直接判断报文是否属于"本端可处理"的范围内，而不需要解析内层IP地址。这大幅简化了VTEP的转发表查找逻辑。

例如，假设某台交换机收到一个VXLAN报文，VNID = 0x010203，其中Region = 0x01、BizID = 0x02、子网 = 0x03。VTEP可以直接将这个VNID映射到对应的VRF（Virtual Routing and Forwarding实例），而无需查看内层IP地址。这意味着，即使两个业务部门使用了完全相同的内网地址段（如10.0.0.0/8），只要BizID不同，VTEP也能正确区分它们的流量。

第二个挑战是性能与成本的平衡。

Overlay封装会引入额外的头部开销。以VXLAN为例，每个报文需要增加50字节的头部：8字节UDP头 + 8字节VXLAN头 + 16字节外层IP/以太网头 + 8字节FCS。这意味着，如果原本的报文是64字节（以太网最小帧），封装后会变成114字节，开销增加了78\%！

对于小包应用来说，这个开销是巨大的。如果一个应用每秒发送100万个64字节的小包，物理网络需要承载114Mbps的带宽，而不是64Mbps。

如何在保持Overlay灵活性的同时，减少封装开销对性能的影响？

阿里云的解决思路是"分层Overlay"。对于性能敏感的业务（如搜索、推荐等高性能计算场景），使用硬件卸载的VXLAN，数据平面完全在网卡上处理。网卡使用专用的硬件逻辑来执行封装/解封装，几乎不占用CPU资源。对于灵活性要求高的业务（如开发测试环境），使用软件Overlay，数据平面在HAVS上处理。软件Overlay的配置更加灵活，支持更多的功能，但性能略低。

这种分层的可行性来自于VXLAN本身的协议设计——VXLAN只是一种封装格式，它不关心封装的"执行者"是软件还是硬件。无论是HAVS还是SmartNIC，只要实现了正确的VXLAN封装/解封装逻辑，就能正确地处理VXLAN流量。这使得两种方案可以无缝共存于同一个VPC网络中。

分层决策的判断标准主要有三个维度：

\textbf{吞吐量需求}：当单台服务器的Overlay吞吐量需求超过10Gbps时，优先考虑硬件卸载方案。因为软件Overlay在处理小包时的CPU消耗是线性增长的，高吞吐量意味着需要占用更多的CPU资源。

\textbf{功能复杂度}：当业务需要动态修改Overlay策略（如频繁的VNI变更、复杂的ACL规则）时，软件Overlay更有优势。硬件Overlay的策略更新通常需要通过固件刷新来实现，周期较长。

\textbf{部署规模}：对于大规模标准化部署（如基础设施层的服务器），硬件Overlay的边际成本更低——网卡上的硬件逻辑可以服务多台虚拟机，不需要为每台虚拟机分配独立的CPU资源。对于小规模灵活部署，软件Overlay更加经济。

这两种方案的共存是通过VPC的控制平面统一管理的。控制平面根据业务需求自动选择合适的Overlay路径，对上层应用透明。这种"软件为主、硬件为辅"的混合策略，平衡了性能和灵活性，是阿里云大规模Overlay网络能够稳定运行的关键。

第三个挑战是运维复杂度的大幅增加。

当Overlay网络覆盖整个集团时，任何一个环节的配置错误都可能导致大范围的网络故障。传统的"telnet到设备上查配置"的方式已经无法适应规模化的要求。

想象一下：当你在telnet到一台核心路由器上敲命令时，可能一不小心敲错了ACL规则，导致整个部门的网络中断。而这种配置错误在大规模网络中可能非常隐蔽——也许只有特定的流量模式才会触发问题，测试阶段可能完全发现不了。

阿里云通过构建"Overlay网络的控制平面"来解决这个问题。

所有的VPC配置、VXLAN隧道建立、MAC地址学习等信息，都通过统一的控制平面进行管理。控制平面运行在专门的服务器集群上，通过BGP、OVSDB等协议与物理网络设备对接。

为什么选择BGP和OVSDB这两种协议？这背后有精心的技术选型考量。

BGP（Border Gateway Protocol）是互联网域间路由的事实标准，它的核心优势是"大规模+高可靠"。BGP天然支持跨域的路由传播，具有丰富的路径属性（如AS-Path、MED、Local-Preference等），可以实现复杂的路由策略。在集团Overlay的场景中，不同业务部门可能属于不同的"自治系统"（AS），BGP可以自然地表达这种归属关系。此外，BGP支持路由聚合和衰减机制，可以有效控制路由表的规模。对于阿里集团这种级别的网络，BGP是唯一能够承载如此大规模路由信息的协议。

OVSDB（Open vSwitch Database）协议则是Open vSwitch生态的核心组成部分，它定义了"控制器与交换机之间的配置同步"机制。OVSDB使用JSON-RPC作为传输协议，数据模型基于RFC 7047定义的OVSDB Schema。通过OVSDB，控制器可以动态地增删VTEP、修改VLAN/VXLAN配置、查询交换机端口状态。这使得网络配置可以"软件定义"——管理员通过API修改配置，由控制器自动同步到所有相关的交换机。

阿里云的控制平面采用了"集中式控制器+BGP/OVSDB数据通道"的混合架构：集中式控制器负责整体策略的决策（如路由计算、安全组规则分发），BGP负责VXLAN隧道端点之间的路由信息传播，OVSDB负责VSwitch层面的配置同步。这种分层设计使得控制平面具有良好的可扩展性——控制器集群可以水平扩展，BGP和OVSDB的分布式特性确保了控制平面本身的可靠性。

当管理员修改VPC配置时——例如增加一条路由、修改一个安全组规则——控制平面会自动将变更同步到所有相关的网络设备。这个同步过程是原子的——要么所有设备都更新成功，要么全部回滚。这避免了部分更新导致的网络不一致。

同时，控制平面还提供了完整的网络拓扑视图和故障诊断工具。当发生网络故障时，管理员可以在控制平面的界面上直观地看到流量路径、检查各节点的配置、定位问题根因。

\textbf{集团Overlay的业务价值}

集团Overlay的全面落地，带来了显著的业务价值。

首先，集团内部的网络架构得到了简化。原来并行的VLAN网络和MPLS网络逐步统一为Overlay网络，减少了运维成本。据估算，每年可节省网络运维人力数百人月。

其次，不同业务部门之间的网络隔离得到了加强。Overlay隧道确保了跨部门通信的安全可控，即使两个部门使用了相同的IP地址，也不会发生串线。

再次，网络弹性大幅提升。新业务上线不再需要申请硬件网络设备，通过Overlay网络可以快速构建隔离的虚拟网络。一台新服务器接入Overlay网络，只需在控制平面注册，即可自动获取网络配置，整个过程可能只需要几分钟。

更重要的是，集团Overlay的实践为阿里云公共云的VPC技术演进提供了宝贵的经验。

许多在集团内部验证过的技术方案——如分层Overlay、热迁移优化、租户地址编码等——随后也被引入到公共云产品中，惠及了更多的云上客户。集团网络和公共云的网络技术在实践中相互促进、共同演进。

\subsection{公网上移：告别硬件NAT}

在云计算的早期，虚拟机访问公网（互联网）通常采用一种叫做"NAT网关"的设备。

NAT网关的作用是将虚拟机发出的私有IP报文，转换（翻译）为公网IP地址，使虚拟机能够访问外部网络；同时，它还要将外部网络返回的报文，转换回对应的私有IP地址，确保通信的双向性。

NAT是一种历史悠久的技术。早期是为了解决IPv4地址不足的问题——多个私有IP设备可以共享一个公网IP访问互联网。在云计算场景下，NAT还有另一个重要作用：隐藏内部网络结构，提高安全性。

在阿里云的Classic网络时代，NAT功能是由专用硬件设备实现的。这些硬件NAT设备被部署在数据中心的关键位置，承担着将大量虚拟机的私网流量转换为公网流量的任务。这种架构在当时是合理的——硬件NAT的性能高、延迟低，能够支撑早期的云计算业务。

\textbf{硬件NAT的局限性}

然而，随着业务规模的增长，硬件NAT的局限性日益凸显。

首先是扩展性问题。硬件NAT设备的容量是固定的——当虚拟机数量超过设备容量时，需要添加更多的NAT设备。这带来了复杂的负载均衡和容错设计。你需要配置NAT设备的高可用（主备或集群模式），需要健康检查来检测设备故障，需要流量调度来均匀分配负载......这些配置和维护工作本身就是一个挑战。

更糟糕的是，当双十一等流量高峰期来临时，NAT设备往往成为瓶颈。假设你购买了1万台ECS实例，它们都需要访问公网。如果NAT设备的总带宽只有100Gbps，当所有实例同时以10Mbps的带宽访问公网时，总带宽需求达到100Gbps，正好达到NAT设备的上限。一旦流量再增加，NAT设备就会成为瓶颈，部分虚拟机无法正常访问公网。

其次是成本问题。高性能硬件NAT设备的价格不菲——一台高端NAT网关的价格可能达到数十万甚至上百万元人民币。随着虚拟机规模的增长，NAT设备的采购和维护成本成为一笔不小的开支。

再次是灵活性问题。硬件NAT设备的配置通常比较僵化，难以适应云环境下快速变化的需求。例如，当需要为某些虚拟机分配固定的公网IP时，硬件NAT往往无法很好地支持。

此外，硬件NAT还带来了单点故障的风险。如果一台NAT设备发生故障，所有依赖它的虚拟机都无法访问公网。虽然可以配置主备来提高可用性，但这增加了复杂度和成本。

\textbf{公网上移的技术架构}

阿里云的"公网上移"方案，就是将NAT功能从专用硬件迁移到通用服务器上，通过软件方式实现。

公网上移这个名字本身就暗示了它的核心理念——不再依赖专用的硬件设备，而是将公网访问的"网关"功能迁移到普通的计算节点上。这些计算节点可以是DPDK-AVS或HAVS，通过软件来实现NAT功能。

公网上移的技术架构有几个关键组件。

首先是软件NAT的实现。阿里云在HAVS/DPDK-AVS上实现了高性能的软件NAT模块，利用DPDK的无锁队列、多核并行等优化技术，软件NAT的性能已经可以媲美甚至超越专用硬件NAT。

软件NAT的工作原理与硬件NAT类似：维护一张NAT表，记录"内网IP:端口"到"公网IP:端口"的映射关系；当私有IP报文离开时，将源IP和端口替换为分配的公网IP和端口；当公网响应返回时，根据NAT表查找对应的私有IP和端口，将报文转发给正确的虚拟机。

测试数据显示，单台DPDK服务器可以支撑数百万的并发NAT会话，吞吐量达到数十Gbps。这个性能已经可以与入门级的硬件NAT设备媲美。

其次是弹性扩缩容能力。由于软件NAT运行在普通服务器上，可以根据负载情况动态调整NAT实例的数量。

当流量增加时，可以快速启动更多的NAT实例；当流量下降时，可以释放部分NAT实例以节省资源。这种弹性能力是硬件NAT无法提供的。

例如，在双十一高峰期，可以临时增加NAT实例来处理突发流量；高峰期过后，这些实例可以被释放，将资源用于其他用途。这种"用多少买多少"的模式，可以大幅节省成本。

再次是统一管控平台。公网上移后，所有的NAT配置（包括端口映射、带宽限制、安全策略等）都通过云平台统一管理。

管理员可以在控制台上配置NAT规则——如某个VPC内的所有虚拟机共享一个公网IP、某个特定的虚拟机独占一个公网IP、某个IP段的所有流量限制在100Mbps等。系统会自动将配置同步到所有NAT实例。

这种集中式的管理方式大大简化了运维工作。以前，管理员需要登录每台NAT设备来配置规则；现在，只需在控制台操作一次，所有设备自动同步。

\textbf{公网上移的额外收益}

公网上移还带来了一个额外的好处：降低了网络架构的复杂度。

在传统的架构中，硬件NAT设备需要与路由器、防火墙、负载均衡器等网络设备配合工作，拓扑复杂、故障点众多。

以一个典型的三层架构为例：虚拟机 $\rightarrow$ 虚拟交换机 $\rightarrow$ 硬件NAT $\rightarrow$ 路由器 $\rightarrow$ 防火墙 $\rightarrow$ 负载均衡器 $\rightarrow$ 互联网。这条路径上有多个设备，任何一个设备故障都可能导致网络中断。

公网上移后，NAT功能集成到了计算节点的虚拟交换机中，网络拓扑变得更加扁平：虚拟机 $\rightarrow$ 虚拟交换机（含NAT） $\rightarrow$ 路由器 $\rightarrow$ 互联网。故障点减少了，排查问题也更容易。

当然，公网上移也面临一些技术挑战。

最大的挑战是性能抖动。当同一台服务器上的多个虚拟机同时发起大量NAT会话时，可能会相互影响。与硬件NAT的专用芯片相比，软件的NAT实现更容易受到邻居虚拟机流量干扰。

例如，当某台虚拟机的流量突然爆发时，可能会占用大量的NAT表项和CPU资源，导致同一台服务器上其他虚拟机的NAT性能下降。这种"吵闹邻居"（Noisy Neighbor）问题在共享基础设施中是常见的。

阿里云通过多种手段来缓解这个问题：CPU资源隔离——NAT处理使用专用的CPU核心，不与其他虚拟机共享；流量调度优化——将NAT流量分散到多个处理线程，避免单线程瓶颈；配额限制——为每个虚拟机设置NAT会话数量的上限。

总的来说，公网上移是阿里云网络虚拟化演进的一个重要里程碑。它证明了软件定义网络的可行性，也为后续的网关产品化（如NAT网关、共享带宽包等）奠定了技术基础。

\section{新时代的VPC：IPv6与跨代迁移}

\subsection{VPC IPv6双栈的工程实践}

2018年，工信部发布了《推进互联网协议第六版（IPv6）规模部署行动计划》，明确要求到2020年底，国内主要互联网服务应完成IPv6改造。作为国内最大的云计算服务商，阿里云面临着艰巨的任务——不仅要让新的VPC实例支持IPv6，还要让存量的IPv4 VPC实例也能平滑接入IPv6网络。

IPv6的好处是众所周知的：近乎无限的地址空间、更简单的报文头部、内置的IPSec支持......但对于已经运行了多年、积累了海量应用和配置的网络来说，向IPv6迁移是一项艰巨的挑战。

在Classic网络向VPC演进的过程中，阿里云积累了大量基于IPv4的实践经验和技术债务。这些经验和技术债务，如何在IPv6时代继续发挥作用，是一个需要深入思考的问题。

\textbf{IPv6带来的新挑战}

从技术角度来看，VPC IPv6面临几个核心挑战。

首先是地址规划的问题。在IPv4时代，阿里云已经为VPC分配了大量的私有地址段。这些地址段是精心规划的，不同的VPC使用不同的地址段，避免冲突。但在IPv6时代，地址空间理论上是无穷无尽的——每个VPC可以分配一个/56甚至/48的IPv6地址段。

这种巨大的地址空间变化，带来了全新的地址规划思路。在IPv4时代，地址是稀缺资源，需要"精打细算"。一个大型企业可能只有几个/24的地址块，需要通过子网划分、VLSM等技术在多个部门之间共享。

在IPv6时代，地址不再是稀缺资源。每个部门、每个应用、甚至每台设备，都可以拥有独立的地址段。地址规划的重点从"如何节省地址"转变为"如何方便管理"。

阿里云在IPv6地址规划中采用了层次化的方案：Region ID + 可用区ID + VPC ID + 虚拟机ID。这种编码方式使得从IPv6地址可以直观地推断出设备的位置和归属，简化了故障排查和安全审计。

以一个具体例子来说明这个编码方案的实际应用。假设阿里云华东1（杭州）Region的代码是"cn-hangzhou"（简化表示），某台虚拟机的IPv6地址为：

\begin{verbatim}
fd00:0000:0c01:0002:0000:0000:0005:0001/128
│   │     │    │   │          │          │
│   │     │    │   │          └─ 虚拟机ID (0005:0001)
│   │     │    │   └─ VPC ID (0000:0000)
│   │     │    └─ 可用区ID (0002)
│   │     └─ Region ID (000c = cn-hangzhou)
└─ 局部唯一地址前缀 (fd00::/8, ULA, 相当于IPv6的"私网地址")
\end{verbatim}

从这个地址我们可以直接读出：这是Region cn-hangzhou、可用区02、VPC ID为0、虚拟机ID为5.1的一台ECS。相比IPv4时代的点分十进制IP地址（172.16.1.5），IPv6的十六进制表示更能直观地反映网络层次关系。

其次是双栈的复杂性。

当一个VPC同时启用IPv4和IPv6时，就进入了"双栈"模式。这意味着：

虚拟机需要同时配置IPv4和IPv6地址。虽然IPv6支持自动地址配置（SLAAC），但在云环境中，通常由平台统一分配地址，这需要相应的逻辑。

网络策略需要同时考虑两种协议。安全组、ACL等规则需要同时支持IPv4和IPv6表达。原来的规则可能需要复制两份，或者需要同时学习两种协议。

网络故障排查的复杂度翻了一倍。同一个连接问题，可能是IPv4的问题，也可能是IPv6的问题，需要分别排查。

再次是DNS的复杂性。在IPv4时代，DNS记录主要是A记录（域名到IPv4地址）。在IPv6时代，需要增加AAAA记录（域名到IPv6地址）。而且，客户端会优先尝试IPv6连接，如果IPv6连接失败再回退到IPv4。这个过程可能导致意外的连接问题。

\textbf{分阶段推进的策略}

阿里云采用了分阶段推进的IPv6策略。

第一阶段，推出NAT64网关。

NAT64是一种将IPv6客户端访问转换为IPv4后端服务的网关产品。对于那些后端服务暂时无法升级到IPv6的场景，NAT64提供了一种过渡方案。

NAT64的工作原理是：网关监听一个IPv6地址，接收来自IPv6客户端的访问请求；将请求中的IPv6目标地址转换为对应的IPv4地址；将转换后的请求转发给后端的IPv4服务器；将服务器返回的IPv4响应转换回IPv6响应，返回给客户端。

对于后端服务来说，整个过程是透明的——它完全不知道客户端使用的是IPv6。

阿里云的NAT64网关部署在VPC的出口网关位置，接收来自外部IPv6网络的访问请求，将其转换为IPv4流量后转发给后端ECS。

这里需要说明一下为什么选择HAVS而非DPDK-AVS来实现NAT64网关。NAT64网关与普通的虚拟交换机不同，它需要维护大规模的会话状态表（用于记录IPv6地址到IPv4地址的映射关系），并且需要动态地进行地址转换。HAVS的Session化架构天然适合这种场景——HAVS维护的会话表（Session Table）可以直接用于NAT64的地址映射查找，无需额外开发。而且，HAVS支持将地址转换的查表操作卸载到网卡硬件（TCAM查找），这使得NAT64网关可以在高性能硬件上实现数十Gbps的转换吞吐量。相比之下，DPDK-AVS虽然性能更高，但它更侧重于"无状态的快速转发"，对于需要维护大量会话状态的NAT64场景，HAVS的统一会话管理能力更为合适。

第二阶段，推出VPC ECS IPv6——让VPC内的虚拟机原生支持IPv6地址。

在这个阶段，虚拟机的网络接口会同时分配一个IPv4地址和一个IPv6地址。用户可以选择只使用IPv6进行通信，或者同时使用双栈。

对于新创建的VPC，默认启用IPv6；对于存量VPC，用户可以手动开启IPv6功能。IPv6地址由平台自动分配，也可以由用户指定。

第三阶段，实现Underlay网络的IPv6升级——这是最困难的一步。

因为Underlay网络涉及大量的物理设备（交换机、路由器等），这些设备的升级周期长、风险高，需要与网络设备厂商紧密配合。

阿里云采取了"软硬分离"的策略：Overlay层（VXLAN隧道、VSwitch等）先行支持IPv6，Underlay层的IPv6升级逐步推进。这意味着，即使物理网络还不支持IPv6，VPC内的虚拟机也可以使用IPv6通信——因为VXLAN封装本身是IP无关的，外层IP可以是IPv4或IPv6。

\textbf{IPv6 ND Proxy的实现}

在双栈实现中，有一个特别值得关注的技术细节：IPv6的ND（Neighbor Discovery）协议与IPv4的ARP协议的不同。

ARP用于解析IPv4地址对应的MAC地址，而ND是IPv6中用于类似目的的协议。ND协议使用ICMPv6报文进行地址解析，比ARP更加复杂，也提供了更多的扩展能力。

在Classic网络中，我们有ARP Proxy来处理ARP请求，防止广播风暴。在VPC中，我们同样需要IPv6 ND Proxy来处理ND请求。

在VPC中实现IPv6支持，需要在虚拟交换机上实现IPv6 ND Proxy功能。这个功能类似于Classic网络中的ARP Proxy——当虚拟机发起ND请求时，ND Proxy截获请求报文，根据本地缓存的ND表项进行响应；如果缓存未命中，则将请求转发到适当的网络位置。

ND Proxy还负责ND表项的维护和超时管理。与ARP Cache类似，ND表项也有生命周期，需要定期刷新。ND Proxy通过发送邻居请求（Neighbor Solicitation）来验证表项的有效性。

\textbf{IPv6的安全考量}

IPv6还带来了一个新的安全问题：SLAAC（Stateless Address Autoconfiguration）攻击。

在SLAAC模式下，IPv6主机可以自行配置地址。攻击者可以发送伪造的路由器公告（Router Advertisement）报文，声称自己是网关，从而让受害者使用攻击者指定的IPv6地址。这样，攻击者就可以实施"中间人"攻击，窃取或篡改通信内容。

阿里云通过在虚拟交换机上实现ND Snooping和RA Guard等安全机制来防范这类攻击。

ND Snooping会监听虚拟机发出的ND报文，记录下"IP $\leftrightarrow$ MAC"的对应关系。如果某个虚拟机声称自己是另一个IP的拥有者，ND Snooping会检测到这个异常并报警或阻止。

RA Guard会过滤掉虚拟机发出的路由器公告报文，防止虚拟机伪装成路由器。在云环境中，只有平台的网络设备才有资格发送RA报文。

阿里云的VPC IPv6实践，为国内云计算行业的IPv6升级提供了宝贵的经验。截至2019年，阿里云已实现了VPC ECS的全面IPv6支持，并持续推进Underlay网络和更多云产品的IPv6改造。

\subsection{Classic迁移到VPC：存量改造的系统工程}

在云计算的世界里，"迁移"是一个永恒的主题。

随着VPC的日益成熟，越来越多的用户开始考虑将业务从Classic网络迁移到VPC。然而，对于已经在线上运行多年、积累了海量资源的用户来说，这个迁移过程绝非易事。

阿里云面临的挑战是：Classic网络中的ECS实例数量以百万计，这些实例承载着不同行业、不同规模客户的业务。将它们迁移到VPC，不仅仅是技术层面的问题，更涉及到用户体验、业务连续性、迁移成本等多方面因素。

\textbf{用户视角的迁移需求}

让我们从用户的视角来看待这个问题。

假设你是一家企业的IT负责人，你的服务器上运行着公司的核心业务系统——电商网站、ERP系统、数据库等。当阿里云告诉你"Classic网络要逐步淘汰了，请迁移到VPC"时，你最关心的是什么？

第一，业务不能中断。公司的业务系统必须7x24小时运行，任何计划外的停机都意味着损失——直接的收入损失、客户信任的下降、员工生产力的降低。

第二，迁移过程要平滑。你不想在半夜被叫起来处理迁移故障，不想在周末加班来应对迁移问题。

第三，可以灰度回滚。如果迁移出了问题，要能够快速回退到之前的状态，而不需要从头开始排查。

第四，迁移进度可控。你想知道迁移需要多久，以便安排后续的工作计划，与业务部门协调时间窗口。

第五，投入产出比合理。迁移本身需要投入人力和时间成本——迁移团队需要熟悉新架构，应用团队需要测试兼容性，运维团队需要学习新工具。你希望这个成本是值得的，而不是"为迁移而迁移"。

这些需求听起来简单，但要全部满足，却需要一套完整的迁移方案和工具链。

\textbf{混挂：最简单的迁移路径}

阿里云为Classic到VPC的迁移提供了几种方案，适用于不同的场景。

首先是"混挂"方案——这是最简单、最直接的迁移路径。

混挂的核心思想是：新创建的ECS实例放在VPC中，而存量ECS实例保持在Classic网络中，通过负载均衡器（SLB）同时挂载两类实例。

这样，来自公网的流量可以同时访问Classic和VPC中的ECS。用户可以逐步将后端服务器从Classic迁移到VPC，而无需改变公网访问的方式。

混挂方案的优势是简单。用户只需要在SLB中添加VPC ECS作为后端，无需修改应用的代码或配置。整个过程可以在线完成，不影响现有服务的可用性。

但它也有局限性。Classic ECS和VPC ECS在同一个SLB后端池中，但它们在内网层面是隔离的。这意味着，Classic ECS无法直接访问VPC ECS的内网服务（如VPC中的RDS、Redis等）。如果应用涉及内网通信，可能需要其他方案。

混挂方案适合的场景是：应用主要通过公网访问，内网通信较少或可以调整。

\textbf{混访：打通云产品的访问}

为了解决内网隔离的问题，阿里云推出了"混访"方案。

云数据库RDS、对象存储OSS等云产品支持同时被Classic网络和VPC网络中的ECS访问。这些产品会提供两个访问域名（经典网络访问域名和VPC访问域名），用户可以根据需要选择使用哪种方式访问。

例如，RDS MySQL可能提供两个访问地址：

经典网络地址：rm-xxx.mysql.rds.aliyuncs.com（解析到经典网络的内网IP）

VPC网络地址：rm-xxx-vpc.mysql.rds.aliyuncs.com（解析到VPC的内网IP）

用户可以根据ECS所在的网络类型，选择相应的访问地址。

当然，随着时间的推移，这个方案也在持续优化和提升用户体验。

混访开启后是经典和VPC两个域名，应用需要根据运行环境选择相应的域名。这点显然用户体验不太友好。那么有没有办法让域名统一呢？

随着PrivateZone（内网DNS）产品的上线，给这个问题的解决带来了契机。通过DNS解析，经典网络和VPC网络相同域名解析到不同的VIP，同时云产品会把DNS包自动路由到正确的网络。

这种方式给用户带来的直接好处是：混访域名统一，ECS上的应用无需做任何修改。应用只需要连接一个域名，DNS会根据ECS所在的网络类型自动返回正确的IP地址。

\textbf{ClassicLink：跨网络的桥梁}

进一步，阿里云还推出了ClassicLink功能，它允许Classic ECS访问VPC内的资源，而无需改变Classic ECS的网络配置。

ClassicLink就像是Classic网络和VPC之间的一座"桥梁"。通过ClassicLink，Classic ECS可以像访问Classic网络内的资源一样，访问VPC内的资源——包括其他VPC内的ECS、RDS、SLB等。

ClassicLink的工作原理是：在Classic网络中注入VPC的路由——当Classic ECS访问VPC内的目标时，报文会被路由到VPC的网关；VPC网关收到报文后，进行正常的VXLAN封装/解封装，将报文送达目标。

ClassicLink使得迁移过程更加平滑。用户可以先通过ClassicLink让Classic ECS访问VPC资源，逐步将应用组件迁移到VPC，最后再完全切换。

\textbf{单ECS迁移：完整的迁移方案}

对于那些需要完全迁移的场景（如公司决定彻底弃用Classic网络），阿里云提供了"单ECS迁移"工具。

这个工具可以帮助用户将单个ECS实例从Classic网络迁移到VPC，迁移过程中保持以下关键属性不变：

公网IP不变。对于依赖固定IP访问的用户来说尤为重要。某些企业可能将IP加入了防火墙白名单，或者在代码中硬编码了服务器的IP地址。如果迁移后IP发生变化，将带来大量的修改工作。

私网IP不变则保证了虚拟机之间内部通信的连续性。由于VPC的Overlay网络支持与Classic网络相同的私网地址段，迁移后的VPC ECS可以使用与原来相同的私网IP地址，无需修改应用程序的网络配置。

迁移过程如下：

第一步，准备阶段。在VPC中创建一个与原Classic ECS相同配置的VPC ECS；在VPC ECS上部署与原ECS相同的应用和数据；配置VPC ECS的网络和安全组，使其可达；进行功能测试，确保应用正常工作。

第二步，切换阶段。通过DNS或负载均衡将流量切换到新的VPC ECS；监控业务指标，确保切换成功。

第三步，清理阶段。确认所有流量已经切换到VPC ECS；释放原Classic ECS。

迁移过程中还有一个关键技术挑战：DNS解析的平滑切换。

在Classic时代，ECS的内网域名解析由阿里云的内网DNS服务器统一管理。迁移到VPC后，用户可以使用PrivateZone自定义DNS解析。PrivateZone允许用户定义自己的DNS记录，覆盖平台默认的解析。

阿里云通过提供迁移工具，帮助用户将原有的DNS记录迁移到PrivateZone。迁移工具会扫描原ECS的配置，提取出DNS解析记录（如/etc/hosts中的条目、应用程序配置的域名等），自动创建对应的PrivateZone记录。

迁移过程中的解析一致性也得到保证。在切换期间，Classic DNS和PrivateZone可以同时生效，确保无论流量切换到哪台服务器，DNS解析都能正确返回。

\textbf{迁移工程的启示}

Classic到VPC的迁移是一项持续多年的系统工程。随着技术的进步和用户需求的演变，阿里云也在不断优化迁移方案和工具。

迁移的本质是"用新架构替换旧架构，同时保持业务连续性"。这个过程不仅仅是技术问题，更是工程管理问题。它需要：

周密的规划——了解迁移的范围、风险、资源需求；
平滑的方案——分阶段迁移、灰度回滚；
高效的执行——自动化工具、标准化的流程；
及时的响应——监控告警、问题排查、应急恢复。

这些原则在未来的网络演进中仍将发挥指导作用。

回过头来看，从Classic到VPC的演进之路，其实折射出了整个云计算网络技术的发展脉络：

早期的"能用就行"——Classic网络虽然有局限性，但满足了快速上线的需求；

中期的"性能优化"——DPDK-AVS、HAVS等解决了性能瓶颈；

后期的"规模化和智能化"——Overlay、硬件卸载、智能运维等支撑了大规模运营。

每一步演进，都是在解决前一步遗留问题的同时，为下一步的发展埋下伏笔。而驱动这个演进的力量，既有业务需求的推动（用户需要更高的性能、更强的隔离），也有技术创新的拉动（新的硬件、新的算法带来了新的可能）。正是这种双向作用，推动着阿里云的网络技术不断向前。

